% Chapter 2, Section 1 _Linear Algebra_ Jim Hefferon
%  http://joshua.smcvt.edu/linearalgebra
%  2001-Jun-10
\chapter{Ruang Vektor}
Pada akhir bab pertama, kita telah cukup memahami cara Metode Gauss
menyelesaikan suatu sistem linear.
Metode tersebut secara sistematis membentuk kombinasi linear dari baris-baris.
Sekarang kita beralih ke kajian umum tentang kombinasi linear.

Kita memerlukan suatu kerangka.
Dalam bab pertama, kadang-kadang kita mengombinasikan vektor-vektor dari
$\Re^2$, kadang-kadang dari $\Re^3$, dan pada kesempatan lain dari ruang
berdimensi lebih tinggi.
Karena itu, dorongan pertama kita mungkin adalah bekerja di $\Re^n$ tanpa
menentukan $n$.
Keuntungannya, setiap hasil akan berlaku sekaligus untuk $\Re^2$, $\Re^3$,
dan banyak ruang lainnya.

Namun, jika penerapan hasil pada banyak ruang sekaligus merupakan suatu
keuntungan, membatasi diri hanya pada ruang-ruang $\Re^n$ justru terlalu
sempit.
Kita ingin agar hasil-hasil tersebut juga berlaku pada kombinasi vektor baris,
seperti dalam bagian terakhir bab pertama.
Kita bahkan telah menjumpai beberapa ruang yang bukan sekadar kumpulan semua
vektor kolom atau semua vektor baris berukuran sama.
Sebagai contoh, kita telah melihat bahwa himpunan solusi suatu sistem homogen
merupakan bidang di dalam $\Re^3$.
Himpunan itu tertutup dalam arti bahwa kombinasi linear dari solusi-solusinya
juga merupakan solusi.
Namun, himpunan tersebut tidak memuat semua vektor kolom bertinggi tiga,
melainkan hanya sebagian di antaranya.

Kita ingin agar hasil tentang kombinasi linear berlaku di mana pun kombinasi
linear dapat didefinisikan dengan masuk akal.
Setiap himpunan semacam itu akan kita sebut \definend{ruang vektor}.
Alih-alih menyatakan hasil kita sebagai ``Setiap kali kita memiliki suatu
kumpulan yang unsur-unsurnya dapat dikombinasikan secara linear dengan masuk
akal \ldots'', kita akan menyatakannya sebagai ``Dalam setiap ruang vektor
\ldots''

Pernyataan seperti itu sekaligus menggambarkan apa yang terjadi dalam banyak
ruang.
Untuk memahami keuntungan beralih dari mempelajari satu ruang ke mempelajari
suatu kelas ruang, perhatikan analogi berikut.
Bayangkan pemerintah membuat undang-undang bagi setiap orang secara terpisah:
``Leslie Jones tidak boleh menyeberang sembarangan.''
Itu tidak baik; pernyataan menjadi lebih ringkas ketika berlaku sekaligus pada
banyak kasus.
Atau, misalkan pemerintah berkata, ``Kim Ke harus berhenti ketika melewati
sebuah kecelakaan.''
Bandingkan dengan, ``Setiap dokter harus berhenti ketika melewati sebuah
kecelakaan.''
Dalam beberapa hal, pernyataan yang lebih umum justru lebih jelas.













\section{Definisi Ruang Vektor}
\index{ruang vektor|(}
Kita akan mempelajari struktur dengan dua operasi, yaitu penjumlahan dan
perkalian skalar, yang memenuhi beberapa syarat sederhana.
Kelak kita akan menelaah syarat-syarat tersebut lebih jauh, tetapi pada
pembacaan pertama perhatikan bahwa semuanya sangat wajar.
Sebagai contoh, setiap operasi yang layak disebut penjumlahan (misalnya
penjumlahan vektor kolom, penjumlahan vektor baris, atau penjumlahan bilangan
real) tentu memenuhi syarat (1) sampai~(5) di bawah ini.



\vspace*{1ex}
\subsection{Definisi dan Contoh}

\begin{definition}
\label{def:VecSpace}
%<*df:VectorSpace>
Suatu \definend{ruang vektor}\index{ruang vektor!definisi}
(di atas \( \Re \)) terdiri atas suatu himpunan \( V \) beserta dua operasi,
`+'\index{vektor!jumlah}\index{jumlah!vektor}\index{penjumlahan vektor}
dan `\( \cdot \)'\index{kelipatan skalar!vektor}\index{vektor!kelipatan skalar},
yang memenuhi syarat-syarat berikut untuk semua vektor
\( \vec{v},\vec{w},\vec{u}\in V \) dan semua
\definend{skalar}\index{skalar} \( r,s\in\Re \):
% \begin{tfae}
\begin{compactenum} 
\item himpunan $V$ tertutup terhadap penjumlahan vektor, yaitu
  \( \vec{v}+\vec{w}\in V \)
\item penjumlahan vektor bersifat komutatif,
  \( \vec{v}+\vec{w}=\vec{w}+\vec{v} \) 
\item penjumlahan vektor bersifat asosiatif,
  \( (\vec{v}+\vec{w})+\vec{u}=\vec{v}+(\vec{w}+\vec{u}) \)
\item terdapat suatu \definend{vektor nol}\index{vektor nol}
    \( \zero\in V \) sedemikian sehingga
    \( \vec{v}+\zero=\vec{v}\, \) untuk semua \( \vec{v}\in V\/ \)
\item setiap \( \vec{v}\in V \) memiliki suatu
    \definend{invers aditif}\index{invers aditif}\index{invers!aditif}
    \( \vec{w}\in V \) sedemikian sehingga \( \vec{w}+\vec{v}=\zero \)
\item himpunan $V$ tertutup terhadap perkalian skalar, yaitu
   \( r\cdot\vec{v}\in V \)
\item perkalian skalar distributif terhadap penjumlahan skalar,
 \( (r+s)\cdot\vec{v}=r\cdot\vec{v}+s\cdot\vec{v} \)
\item perkalian skalar distributif terhadap penjumlahan vektor,
  \( r\cdot(\vec{v}+\vec{w})=r\cdot\vec{v}+r\cdot\vec{w} \)
\item perkalian biasa antarskalar bersifat asosiatif dengan perkalian skalar,
  \( (rs)\cdot\vec{v} =r\cdot(s\cdot\vec{v}) \)
\item perkalian dengan skalar~$1$ merupakan operasi identitas,
  \( 1\cdot\vec{v}=\vec{v} \).
\end{compactenum}
% \end{tfae}  
%</df:VectorSpace>
\end{definition}

\vspace{-4.5ex}  % why is the extra space there?  Bug in shading code?
\begin{remark}
Definisi tersebut melibatkan dua jenis penjumlahan dan dua jenis perkalian,
sehingga pada awalnya mungkin tampak membingungkan.
Sebagai contoh, dalam syarat~(7), `$+$' di ruas kiri adalah penjumlahan dua
bilangan real, sedangkan `$+$' di ruas kanan adalah penjumlahan dua vektor
dalam \( V\/ \).
Ungkapan-ungkapan tersebut tidak ambigu karena konteksnya; misalnya,
\( r \) dan \( s \) adalah bilangan real sehingga `\( r+s \)' hanya mungkin
berarti penjumlahan bilangan real.
Demikian pula, `$rs$' di ruas kiri butir~(9) adalah perkalian biasa bilangan
real, sedangkan `$s\cdot\vec{v}$' di ruas kanannya adalah perkalian skalar
yang didefinisikan untuk ruang vektor ini.
\end{remark}

\vspace{-0.5ex}  % why is the extra space there?  Bug in shading code?
Cara terbaik memahami definisi tersebut adalah menelaah contoh-contoh di
bawah ini dan memeriksa kesepuluh syarat pada masing-masing contoh.
Contoh pertama memuat pemeriksaan itu secara terperinci.
Gunakan contoh tersebut sebagai pola bagi contoh-contoh lainnya.
Syarat \definend{ketertutupan}\index{ruang vektor!ketertutupan}, yaitu
(1) dan~(6), sangatlah penting.
Keduanya menyatakan bahwa operasi penjumlahan dan perkalian skalar selalu
bermakna\Dash operasi tersebut didefinisikan untuk setiap pasangan vektor
dan untuk setiap pasangan skalar dan vektor, serta hasil operasinya merupakan
anggota himpunan itu.


\begin{example}  \label{PlaneThruOriginSubsp}
Subhimpunan \( \Re^2 \) berikut merupakan garis yang melalui titik asal.
\begin{equation*}
  L=\set{ \colvec{x \\ y}  \suchthat y=3x}
\end{equation*}
Kita akan memverifikasi bahwa himpunan tersebut merupakan ruang vektor dengan
makna biasa dari `+' dan `\(\cdot\)'.
\begin{equation*}
  \colvec{x_1 \\ y_1}
  +
  \colvec{x_2 \\ y_2}
  =
  \colvec{x_1+x_2 \\ y_1+y_2}
  \qquad
  r\cdot
  \colvec{x \\ y}
  =
  \colvec{rx \\ ry}
\end{equation*}
Operasi-operasi ini hanyalah operasi biasa yang digunakan kembali pada
subhimpunan $L$.
Kita mengatakan bahwa \( L \) \definend{mewarisi\/}\index{operasi yang diwarisi}
operasi-operasi tersebut dari \( \Re^2 \).

Kita akan memeriksa kesepuluh syarat.
Lima syarat pertama berkaitan dengan penjumlahan.
Untuk syarat~(1), yakni ketertutupan terhadap penjumlahan, misalkan kita
memulai dengan dua vektor dari garis~$L$,
\begin{equation*}
  \vec{v}_1=\colvec{x_1 \\ y_1}
  \quad
  \vec{v}_2=\colvec{x_2 \\ y_2}
\end{equation*}
sehingga keduanya memenuhi syarat $y_1=3x_1$ dan~$y_2=3x_2$.
Jumlah keduanya,
\begin{equation*}
  \vec{v}_1+\vec{v}_2
  =\colvec{x_1+x_2 \\ y_1+y_2}
\end{equation*}
juga merupakan anggota garis $L$ karena komponen keduanya tiga kali komponen
pertamanya: $y_1+y_2=3(x_1+x_2)$ mengikuti syarat pada $\vec{v}_1$
dan~$\vec{v}_2$.
Untuk~(2), yakni komutativitas penjumlahan vektor, bandingkan
\begin{equation*}
  \vec{v}_1+\vec{v}_2
  =\colvec{x_1+x_2 \\ y_1+y_2}
  \quad
  \vec{v}_2+\vec{v}_1
  =\colvec{x_2+x_1 \\ y_2+y_1}
\end{equation*}
dan perhatikan bahwa keduanya sama karena entri-entrinya merupakan bilangan
real dan penjumlahan bilangan real bersifat komutatif.
(Syarat bahwa vektor-vektor tersebut terletak pada garis tidak relevan bagi
syarat ini; keduanya dapat dipertukarkan hanya karena penjumlahan semua vektor
di bidang bersifat komutatif.)
Syarat~(3), asosiativitas penjumlahan vektor, serupa.
\begin{align*}
  (\colvec{x_1 \\ y_1}
  +\colvec{x_2 \\ y_2})
  +\colvec{x_3 \\ y_3}
  &=\colvec{(x_1+x_2)+x_3 \\ (y_1+y_2)+y_3}  \\
  &=\colvec{x_1+(x_2+x_3) \\ y_1+(y_2+y_3)}  \\
  &=\colvec{x_1 \\ y_1}
  +(\colvec{x_2 \\ y_2}
  +\colvec{x_3 \\ y_3})
\end{align*}
Untuk syarat keempat, kita harus memberikan suatu vektor yang berperan sebagai
unsur nol. Vektor yang semua entrinya nol memenuhi keperluan tersebut.
\begin{equation*}
  \colvec{x \\ y}
  +\colvec{0 \\ 0}
  =\colvec{x \\ y}
\end{equation*}
Perhatikan bahwa $\zero\in L$ karena komponen keduanya tiga kali komponen
pertamanya.
Untuk~(5), yakni bahwa bagi setiap~$\vec{v}\in L$ kita dapat memberikan suatu
invers aditif, kita memiliki
\begin{equation*}
  \colvec{-x \\ -y}
  +\colvec{x \\ y}
  =\colvec{0 \\ 0}
\end{equation*}
sehingga vektor $-\vec{v}$ merupakan invers yang dicari.
Seperti pada syarat sebelumnya, perhatikan bahwa jika $\vec{v}\in L$, sehingga
$y=3x$, maka $-\vec{v}\in L$ juga karena $-y=3(-x)$.

Pemeriksaan kelima syarat yang berkaitan dengan perkalian skalar serupa.
Untuk~(6), yakni ketertutupan terhadap perkalian skalar, misalkan $r\in\R$
dan $\vec{v}\in L$, yaitu
\begin{equation*}
  \vec{v}=\colvec{x \\ y}
\end{equation*}
memenuhi~$y=3x$. Maka
\begin{equation*}
  r\cdot\vec{v}=r\cdot\colvec{x \\ y}
  =\colvec{rx \\ ry}
\end{equation*}
juga merupakan anggota $L$:~hubungan $ry=3\cdot rx$ berlaku karena $y=3x$.
Selanjutnya, persamaan berikut memeriksa syarat~(7).
\begin{equation*}
  (r+s)\cdot\colvec{x \\ y}
  =\colvec{(r+s)x \\ (r+s)y}
  =\colvec{rx+sx \\ ry+sy}
  =r\cdot\colvec{x \\ y}+s\cdot\colvec{x \\ y}
\end{equation*}
Untuk~(8), kita memperoleh
\begin{equation*}
  r\cdot(\colvec{x_1 \\ y_1}+\colvec{x_2 \\ y_2})
  =\colvec{r(x_1+x_2) \\ r(y_1+y_2)}
  =\colvec{rx_1+rx_2 \\ ry_1+ry_2}
  =r\cdot\colvec{x_1 \\ y_1}+r\cdot\colvec{x_2 \\ y_2}
\end{equation*}
Syarat kesembilan,
\begin{equation*}
  (rs)\cdot\colvec{x \\ y}
  =\colvec{(rs)x \\ (rs)y}
  =\colvec{r(sx) \\ r(sy)}
  =r\cdot(s\cdot\colvec{x \\ y})
\end{equation*}
dan syarat kesepuluh juga langsung diperoleh.
\begin{equation*}
  1\cdot\colvec{x \\ y}
  =\colvec{1x \\ 1y}
  =\colvec{x \\ y}
\end{equation*}
\end{example}

\begin{example}  \label{ex:RealVecSpaces}
Seluruh bidang, yakni himpunan \( \Re^2 \), merupakan ruang vektor dengan
operasi `\( + \)' dan `\( \cdot \)' bermakna seperti biasa.
\begin{equation*}
  \colvec{x_1 \\ y_1}
  +
  \colvec{x_2 \\ y_2}
  =
  \colvec{x_1+x_2 \\ y_1+y_2}
  \qquad
  r\cdot
  \colvec{x \\ y}
  =
  \colvec{rx \\ ry}
\end{equation*}
Kita hanya akan memeriksa dua syarat, yaitu syarat ketertutupan.

Untuk (1), perhatikan bahwa hasil penjumlahan vektor
\begin{equation*}
  \colvec{x_1 \\ y_1}
  +\colvec{x_2 \\ y_2}
  =\colvec{x_1+x_2 \\ y_1+y_2}
\end{equation*}
merupakan susunan kolom dengan dua entri real sehingga menjadi anggota
bidang~\( \Re^2 \).
Berbeda dengan contoh sebelumnya, di sini tidak ada syarat yang membatasi
komponen pertama dan kedua vektor-vektornya.

Syarat (6) serupa. Vektor
\begin{equation*}
  r\cdot\colvec{x \\ y}
  =\colvec{rx \\ ry}
\end{equation*}
memiliki dua entri real sehingga menjadi anggota~\( \Re^2 \).
\end{example}

Dengan cara yang serupa, setiap \( \Re^n \) merupakan ruang vektor dengan
operasi penjumlahan vektor dan perkalian skalar yang biasa.
(Dalam \( \Re^1 \), kita biasanya tidak menuliskan anggotanya sebagai vektor
kolom; jadi, kita biasanya tidak menulis `\( (\pi) \)', melainkan cukup
`\( \pi \)'.)

\begin{example} \label{ex:ColsIntEntNotVS}
\nearbyexample{PlaneThruOriginSubsp} memberikan subhimpunan $\Re^2$ yang
merupakan ruang vektor.
% \nearbyexample{ex:RealVecSpaces} shows that the entire
% set of two-tall vectors with real entries is also a vector space. 
Sebagai pembanding, perhatikan himpunan kolom bertinggi dua dengan entri-entri
bilangan bulat, di bawah operasi penjumlahan komponen demi komponen dan
perkalian skalar yang sama.
Himpunan ini merupakan subhimpunan $\Re^2$, tetapi bukan ruang vektor: himpunan
tersebut tidak tertutup terhadap perkalian skalar, sehingga tidak memenuhi
syarat~(6).
Sebagai contoh, di ruas kiri berikut terdapat sebuah skalar dan sebuah vektor
yang entri-entrinya bilangan bulat.
\begin{equation*}
  0.5
  \cdot
  \colvec[r]{4 \\ 3}
  =
  \colvec[r]{2 \\ 1.5}
\end{equation*}
Di ruas kanan terdapat vektor kolom yang bukan anggota himpunan tersebut karena
tidak semua entrinya merupakan bilangan bulat.
\end{example}

\begin{example}  \label{ex:TrivSbspReFour}
Himpunan beranggota tunggal
\begin{equation*}
  \set{ \colvec[r]{0 \\ 0 \\ 0 \\ 0} }
\end{equation*}
merupakan ruang vektor di bawah operasi
\begin{equation*}
  \colvec[r]{0 \\ 0 \\ 0 \\ 0}
  +
  \colvec[r]{0 \\ 0 \\ 0 \\ 0}
  =
  \colvec[r]{0 \\ 0 \\ 0 \\ 0}
  \qquad
  r\cdot
  \colvec[r]{0 \\ 0 \\ 0 \\ 0}
  =
  \colvec[r]{0 \\ 0 \\ 0 \\ 0}
\end{equation*}
yang diwarisinya dari \( \Re^4 \).
\end{example}

Suatu ruang vektor harus memiliki sedikitnya satu anggota, yaitu vektor
nolnya. Karena itu, ruang vektor beranggota tunggal merupakan ruang vektor
terkecil yang mungkin.

\begin{definition} \label{df:TrivialVectorSpace}
%<*df:TrivialVectorSpace>
Ruang vektor beranggota tunggal disebut ruang
\definend{trivial}\index{ruang vektor!trivial}%
\index{ruang trivial}.
%</df:TrivialVectorSpace>
\end{definition}

Contoh-contoh sejauh ini melibatkan himpunan vektor kolom dengan operasi yang
biasa. Namun, ruang vektor tidak harus berupa kumpulan vektor kolom, bahkan
tidak harus berupa kumpulan vektor baris.
Di bawah ini terdapat beberapa jenis ruang vektor lainnya.
Istilah `ruang vektor' tidak berarti `kumpulan kolom bilangan real'.
Maknanya lebih dekat dengan `kumpulan tempat setiap kombinasi linear dapat
didefinisikan dengan masuk akal'.

\begin{example} \label{ex:PolySpaceThree}
Perhatikan
\( \polyspace_3=\set{a_0+a_1x+a_2x^2+a_3x^3\suchthat a_0,\ldots,a_3\in\Re} \),
yaitu himpunan polinomial berderajat tiga atau kurang
(dalam buku ini, polinomial konstan, termasuk polinomial nol, kita anggap
berderajat nol).
Himpunan tersebut merupakan ruang vektor di bawah operasi
\begin{multline*}
   (a_0+a_1x+a_2x^2+a_3x^3)+(b_0+b_1x+b_2x^2+b_3x^3)  \\
     =(a_0+b_0)+(a_1+b_1)x+(a_2+b_2)x^2+(a_3+b_3)x^3
\end{multline*}
dan
\begin{equation*}
   r\cdot(a_0+a_1x+a_2x^2+a_3x^3)=
     (ra_0)+(ra_1)x+(ra_2)x^2+(ra_3)x^3
\end{equation*}
(verifikasinya mudah).
Ruang vektor ini patut diperhatikan karena operasi-operasi tersebut merupakan
operasi polinomial yang dikenal dari aljabar sekolah menengah.
Sebagai contoh,
$
  3\cdot(1-2x+3x^2-4x^3)-2\cdot(2-3x+x^2-(1/2)x^3)=-1+7x^2-11x^3$.

Walaupun ruang ini bukan subhimpunan dari suatu \( \Re^n \), dalam pengertian
tertentu kita dapat memandang $\polyspace_3$ sebagai sesuatu yang ``sama''
dengan \( \Re^4 \).
Jika kita mengidentifikasi anggota kedua ruang itu dengan cara berikut,
\begin{equation*}
  a_0+a_1x+a_2x^2+a_3x^3
  \quad\text{berkorespondensi dengan}\quad
  \colvec{a_0 \\ a_1 \\ a_2 \\ a_3}
\end{equation*}
maka operasi-operasinya juga saling berkorespondensi.
Berikut adalah contoh penjumlahan yang saling berkorespondensi.
\begin{equation*}
  \mbox{\begin{tabular}{lr}
       &\(1-2x+0x^2+1x^3\) \\
     + &\(2+3x+7x^2-4x^3\) \\ \hline
       &\(3+1x+7x^2-3x^3\)
  \end{tabular}  }
  \quad\text{berkorespondensi dengan}\quad
  \colvec[r]{1 \\ -2 \\ 0 \\ 1}
  +
  \colvec[r]{2 \\ 3 \\ 7 \\ -4}
  =
  \colvec[r]{3 \\ 1 \\ 7 \\ -3}
\end{equation*}
Objek-objek yang kita pandang sebagai sesuatu yang ``sama'' dijumlahkan menjadi
hasil yang ``sama''.
Bab Tiga akan menjadikan gagasan korespondensi ruang vektor ini lebih tepat.
Untuk saat ini, kita cukup mempertahankannya sebagai suatu intuisi.
\end{example}

Secara umum, kita menulis \( \polyspace_n \) untuk ruang vektor polinomial
berderajat~$n$ atau kurang,
\( \set{a_0+a_1x+a_2x^2+\cdots+a_nx^n\suchthat a_0,\ldots,a_n\in\Re} \),
di bawah operasi penjumlahan polinomial dan perkalian skalar yang biasa.
\index{ruang vektor!polinomial}
Kita akan sering menggunakan ruang-ruang ini sebagai contoh.

\begin{example} \label{ex:MatSpaceTwoByTwo}
Himpunan \( \matspace_{\nbym{2}{2}} \) yang terdiri atas matriks-matriks
\( \nbym{2}{2} \) dengan entri bilangan real merupakan ruang vektor di bawah
operasi alami entri demi entri.
\begin{equation*}
  \begin{mat}
    a  &b \\
    c  &d
  \end{mat}
  +
  \begin{mat}
    w  &x \\
    y  &z
  \end{mat}
  =
  \begin{mat}
    a+w  &b+x \\
    c+y  &d+z
  \end{mat}
  \qquad
  r\cdot
  \begin{mat}
    a  &b \\
    c  &d
  \end{mat}
  =
  \begin{mat}
    ra  &rb \\
    rc  &rd
  \end{mat}
\end{equation*}
Seperti dalam contoh sebelumnya, kita dapat memandang ruang ini sebagai
sesuatu yang ``sama'' dengan \( \Re^4 \).
\end{example}

Kita menulis \( \matspace_{\nbym{n}{m}} \) untuk ruang vektor
matriks-matriks $\nbym{n}{m}$ di bawah operasi alami penjumlahan matriks dan
perkalian skalar.\index{ruang vektor!matriks}
Seperti ruang polinomial, ruang-ruang ini akan sering kita gunakan sebagai
contoh.

\begin{example}   \label{ex:FcnsNToRIsVecSp}
Himpunan \( \set{f\suchthat \map{f}{\N}{\Re} } \) yang terdiri atas semua
fungsi bernilai real dari satu variabel bilangan asli merupakan ruang vektor
di bawah operasi
\begin{equation*}
  (f_1+f_2)\,(n)=f_1(n)+f_2(n)
  \qquad
  (r\cdot f)\,(n)=r\,f(n)
\end{equation*}
sehingga, sebagai contoh, jika \( f_1(n)=n^2+2\sin(n) \) dan
\( f_2(n)=-\sin(n)+0.5 \), maka \( (f_1+2f_2)\,(n)=n^2+1 \).

Kita dapat memandang ruang ini sebagai perumuman dari
\nearbyexample{ex:RealVecSpaces}\Dash alih-alih vektor bertinggi~$2$,
fungsi-fungsi ini menyerupai vektor bertinggi tak hingga.
\begin{equation*}
  \text{
    \begin{tabular}{c|c}
      \( n \)      &\( f(n)=n^2+1 \)  \\ \hline
      \( 0      \) &\( 1      \)    \\
      \( 1      \) &\( 2      \)    \\
      \( 2      \) &\( 5      \)    \\
      $3$          &$10$            \\
      \( \vdots \) &\( \vdots \)    
    \end{tabular} }
    \quad\text{berkorespondensi dengan}\quad
    \colvec[r]{
           1           \\
           2           \\
           5           \\
           10          \\
           \vdotswithin{10}      }
\end{equation*}
Penjumlahan dan perkalian skalar dilakukan komponen demi komponen, seperti
dalam \nearbyexample{ex:RealVecSpaces}.
(Kita dapat memformalkan istilah ``bertinggi tak hingga'' dengan mengatakan
bahwa objek tersebut merupakan barisan tak hingga, atau fungsi dari $\N$
ke $\Re$.)
\end{example}

\begin{example} \label{ex:PolysOfAllFiniteDegrees}
Himpunan polinomial dengan koefisien real
\begin{equation*}
 \set{ a_0+a_1x+\cdots+a_nx^n\suchthat n\in\N
    \text{ dan } a_0,\ldots,a_n\in\Re}
\end{equation*}
menjadi ruang vektor jika dilengkapi dengan operasi alami `$+$',
\begin{multline*}
  (a_0+a_1x+\cdots+a_nx^n)+(b_0+b_1x+\cdots+b_nx^n)  \\
     =(a_0+b_0)+(a_1+b_1)x+\cdots +(a_n+b_n)x^n
\end{multline*}
dan `$\cdot$'.
\begin{equation*}
  r\cdot (a_0+a_1x+\ldots a_nx^n)
   =
  (ra_0)+(ra_1)x+\ldots (ra_n)x^n
\end{equation*}
Ruang ini berbeda dari ruang $\polyspace_3$ dalam
\nearbyexample{ex:PolySpaceThree}.
Ruang ini tidak hanya memuat polinomial berderajat~tiga, tetapi juga
polinomial berderajat~tiga puluh dan berderajat~tiga ratus.
Setiap polinomial tentu berderajat hingga, tetapi tidak ada satu batas yang
berlaku bagi derajat semua anggota himpunan tersebut.

Seperti pada contoh sebelumnya, kita dapat memikirkan contoh ini melalui tupel
tak hingga.
Sebagai contoh, kita dapat memandang \( 1+3x+5x^2 \) sebagai sesuatu yang
berkorespondensi dengan \( (1,3,5,0,0,\ldots) \).
Namun, ruang ini berbeda dari ruang dalam
\nearbyexample{ex:FcnsNToRIsVecSp}.
Di sini, setiap anggota himpunan berderajat hingga; dengan kata lain, di bawah
korespondensi tersebut tidak ada anggota ruang ini yang cocok dengan
\( (1,2,5,10,\,\ldots\,) \).
Vektor-vektor dalam ruang ini berkorespondensi dengan tupel tak hingga yang
akhirnya hanya berisi nol.
\end{example}

\begin{example}  \label{ex:RealValuedFcns}
Himpunan
\( \set{f\suchthat \map{f}{\Re}{\Re} } \)
yang terdiri atas semua fungsi bernilai real dari satu variabel real merupakan
ruang vektor di bawah operasi berikut.
\begin{equation*}
  (f_1+f_2)\,(x)=f_1(x)+f_2(x)
  \qquad
  (r\cdot f)\,(x)=r\,f(x)
\end{equation*}
Perbedaan antara ruang ini dan \nearbyexample{ex:FcnsNToRIsVecSp} terletak
pada domain fungsi-fungsinya.
\end{example}

\begin{example}\label{ex:ACos+BSin}
Himpunan
\( F=\{ a\cos\theta+b\sin\theta \suchthat a,b\in\Re\} \)
yang terdiri atas fungsi-fungsi bernilai real dari variabel real \( \theta \)
merupakan ruang vektor di bawah operasi
%\nearbyexample{ex:RealValuedFcns}:
\begin{equation*}
  (a_1\cos\theta+b_1\sin\theta)+(a_2\cos\theta+b_2\sin\theta)
    =(a_1+a_2)\cos\theta+(b_1+b_2)\sin\theta
\end{equation*}
dan
\begin{equation*}
  r\cdot (a\cos\theta+b\sin\theta)
   =(ra)\cos\theta+(rb)\sin\theta
\end{equation*}
yang diwarisi dari ruang dalam contoh sebelumnya.
(Kita dapat memandang \( F \) sebagai sesuatu yang ``sama'' dengan \( \Re^2 \)
karena $a\cos\theta+b\sin\theta$ berkorespondensi dengan vektor yang
komponen-komponennya $a$ dan $b$.)
\end{example}

\begin{example}  \label{ex:SpaceSatisfyingDiffQ}
Himpunan
\begin{equation*}
  \set{\map{f}{\Re}{\Re}\suchthat \dfrac{d^2f}{dx^2}+f=0}
\end{equation*}
merupakan ruang vektor di bawah penafsiran yang kini sudah terasa alami.
\begin{equation*}
  (f+g)\,(x)=f(x)+g(x)
  \qquad
  (r\cdot f)\,(x)=r\,f(x)
\end{equation*}
Secara khusus, perhatikan bahwa Kalkulus dasar memberikan
\begin{equation*}
   \frac{d^2(f+g)}{dx^2}+(f+g)
   =(\frac{d^2f}{dx^2}+f)+(\frac{d^2g}{dx^2}+g)
\end{equation*}
dan
\begin{equation*}
   \frac{d^2(rf)}{dx^2}+(rf)
   =r(\frac{d^2 f}{dx^2}+f)
\end{equation*}
sehingga ruang tersebut tertutup terhadap penjumlahan dan perkalian skalar.
Ruang ini ternyata sama dengan ruang dalam contoh sebelumnya\Dash fungsi yang
memenuhi persamaan diferensial ini berbentuk
$a\cos\theta+b\sin\theta$\Dash tetapi deskripsi ini juga mengisyaratkan
perluasan ke himpunan solusi persamaan diferensial lainnya.
\end{example}

\begin{example} \label{ex:HomoSlnMakesVS}
Himpunan solusi suatu sistem linear homogen dalam \( n \) variabel merupakan
ruang vektor di bawah operasi yang diwarisi dari \( \Re^n \).
Sebagai contoh, untuk memeriksa ketertutupan terhadap penjumlahan, perhatikan
suatu persamaan biasa dalam sistem tersebut,
$c_1x_1+\cdots+c_nx_n=0$, dan misalkan kedua vektor berikut
\begin{equation*}
   \vec{v}=\colvec{v_1 \\ \vdotswithin{v_1} \\ v_n}
   \qquad
   \vec{w}=\colvec{w_1 \\ \vdotswithin{w_1} \\ w_n}
\end{equation*}
memenuhi persamaan itu.
Maka jumlah keduanya, \( \vec{v}+\vec{w}\/ \), juga memenuhi persamaan tersebut:
\(
  c_1(v_1+w_1)+\cdots+c_n(v_n+w_n)
  =(c_1v_1+\cdots+c_nv_n)+(c_1w_1+\cdots+c_nw_n)
  =0
\).
Pemeriksaan syarat ruang vektor lainnya juga langsung dilakukan.
\end{example}

Kita sering menghilangkan lambang perkalian `\( \cdot \)' di antara skalar
dan vektor.
Kita membedakan perkalian dalam \( c_1v_1 \) dari perkalian dalam
\( r\vec{v}\, \) berdasarkan konteks: jika kedua faktor merupakan bilangan
real, operasinya tentu perkalian antarbilangan real; jika salah satunya vektor,
operasinya tentu perkalian skalar dengan vektor.

\nearbyexample{ex:HomoSlnMakesVS} membawa kita kembali ke titik awal karena
contoh itu merupakan salah satu motivasi kita.
Kini, setelah memperoleh gambaran tentang jenis struktur yang memenuhi
definisi ruang vektor, kita dapat menelaah kembali definisi tersebut.
Sebagai contoh, mengapa definisi itu secara khusus memuat syarat
\( 1\cdot\vec{v}=\vec{v} \), tetapi tidak memuat syarat
\( 0\cdot\vec{v}=\zero \)?

Salah satu jawabannya adalah bahwa ini memang sebuah definisi\Dash definisi
tersebut memberikan aturan yang harus kita ikuti untuk melanjutkan.

Jawaban lain mungkin lebih memuaskan.
Para ahli di bidang ini telah berupaya menemukan keseimbangan yang tepat antara
kekuatan dan keumuman.
Definisi ini disusun agar memuat syarat-syarat yang diperlukan untuk
membuktikan semua sifat menarik dan penting dari ruang kombinasi linear.
Dalam pembahasan selanjutnya, kita akan menurunkan semua sifat yang wajar bagi
kumpulan kombinasi linear dari syarat-syarat dalam definisi tersebut.

Hasil berikut merupakan suatu contoh.
Kita tidak perlu memasukkan sifat-sifat ini ke dalam definisi ruang vektor
karena semuanya mengikuti sifat-sifat yang sudah tercantum di sana.

\begin{lemma} \label{lm:ElementaryPropertiesOfVectorSpaces}
%<*lm:ElementaryPropertiesOfVectorSpaces>
Dalam setiap ruang vektor \( V \), untuk setiap \( \vec{v}\in V \) dan
\( r\in\Re \), berlaku
(1)~\( 0\cdot\vec{v}=\zero \),
(2)~\( (-1\cdot\vec{v})+\vec{v}=\zero \), dan
(3)~\( r\cdot\zero=\zero \).
%</lm:ElementaryPropertiesOfVectorSpaces>
\end{lemma}

\begin{proof}
%<*pf:ElementaryPropertiesOfVectorSpaces0>
Untuk (1), perhatikan bahwa
\( \vec{v}=(1+0)\cdot\vec{v}=\vec{v}+(0\cdot\vec{v}) \).
Tambahkan invers aditif dari \( \vec{v} \) pada kedua ruas, yaitu vektor
\( \vec{w} \) sedemikian sehingga \( \vec{w}+\vec{v}=\zero \).
\begin{align*}
  \vec{w}+\vec{v}
  &=\vec{w}+\vec{v}+0\cdot\vec{v}  \\
  \zero
  &=\zero+0\cdot\vec{v}                   \\
  \zero
  &=0\cdot\vec{v}
\end{align*}
Butir~(2) mudah:
\(  (-1\cdot\vec{v})+\vec{v}=(-1+1)\cdot\vec{v}=0\cdot\vec{v}=\zero \).
Untuk (3),
\( r\cdot\zero
  =
  r\cdot(0\cdot\zero)
  =
  (r\cdot 0)\cdot\zero
  =
  \zero \)
memberikan hasil yang diinginkan.
%</pf:ElementaryPropertiesOfVectorSpaces0>
\end{proof}

Butir kedua menunjukkan bahwa kita dapat menulis invers aditif dari
\( \vec{v} \) sebagai `\( -\vec{v}\, \)' tanpa khawatir akan kebingungan
dengan \( (-1)\cdot\vec{v} \).

\medskip
Sebagai rangkuman, kajian reduksi Gauss dalam Bab Satu membawa kita untuk
meninjau kumpulan kombinasi linear.
Karena itu, dalam bab ini kita mendefinisikan ruang vektor sebagai struktur
tempat kita dapat membentuk kombinasi semacam itu,
% expressions of the form \( c_1\cdot\vec{v}_1+\dots+c_n\cdot\vec{v}_n \)
dengan beberapa syarat sederhana pada operasi penjumlahan dan perkalian
skalar. Singkatnya:~ruang vektor merupakan konteks yang tepat untuk
mempelajari linearitas.

% Finally, a comment.
Karena sistem linear mengisi satu bab penuh, terlebih lagi bab pertama,
pembaca mungkin menduga bahwa tujuan buku ini adalah mempelajari sistem
linear. Kenyataannya, kita bukan terutama akan menggunakan ruang vektor untuk
mempelajari sistem linear; justru sistem linear menjadi titik awal untuk
mempelajari ruang vektor.
Beragam contoh dalam subbagian ini menunjukkan bahwa kajian ruang vektor
menarik dan penting dengan sendirinya.
Sistem linear tidak akan menghilang dari pembahasan.
Namun, mulai sekarang objek kajian utama kita adalah ruang vektor.

\begin{exercises}
  \item 
    Sebutkan vektor nol untuk setiap ruang vektor berikut.
    \begin{exparts}
      \partsitem Ruang polinomial berderajat tiga di bawah operasi alaminya.
      \partsitem Ruang matriks \( \nbym{2}{4} \).
      \partsitem Ruang
        \( \set{\map{f}{[0..1]}{\Re}\suchthat f\text{ kontinu}} \).
      \partsitem Ruang fungsi bernilai real dari satu variabel bilangan asli.
    \end{exparts}
    \begin{answer}
      \begin{exparts}
        \partsitem \( 0+0x+0x^2+0x^3 \)
        \partsitem \( \begin{mat}[r]
                   0  &0  &0  &0  \\
                   0  &0  &0  &0
                 \end{mat} \)
        \partsitem Fungsi konstan \( f(x)=0 \)
        \partsitem Fungsi konstan \( f(n)=0 \)
      \end{exparts}  
    \end{answer}
  \recommended \item
    Tentukan invers aditif dari vektor berikut di dalam ruang vektornya.
    \begin{exparts}
      \partsitem Dalam \( \polyspace_3 \), vektor \( -3-2x+x^2 \).
      \partsitem Dalam ruang \( \nbyn{2} \),
        \begin{equation*}
          \begin{mat}[r]
            1  &-1  \\
            0  &3
          \end{mat}.
        \end{equation*}
     \partsitem Dalam \( \set{ae^x+be^{-x}\suchthat a,b\in\Re} \), yaitu
       ruang fungsi dari variabel real \( x \) di bawah operasi alaminya,
       vektor \( 3e^x-2e^{-x} \).
    \end{exparts}
    \begin{answer}
      \begin{exparts*}
        \partsitem \( 3+2x-x^2 \)
        \partsitem \( \begin{mat}[r]
                   -1  &+1  \\
                    0  &-3
                 \end{mat} \)
        \partsitem \( -3e^x+2e^{-x} \)
      \end{exparts*}  
    \end{answer}
  \recommended \item 
    Untuk masing-masing himpunan, daftarkan tiga anggotanya lalu tunjukkan
    bahwa himpunan tersebut merupakan ruang vektor.
    \begin{exparts}
      \partsitem Himpunan polinomial linear
        \( \polyspace_1=\set{a_0+a_1x\suchthat a_0,a_1\in\Re} \) di bawah
        operasi penjumlahan polinomial dan perkalian skalar yang biasa.
      \partsitem Himpunan polinomial linear
        \( \set{a_0+a_1x\suchthat a_0-2a_1=0} \), di bawah operasi
        penjumlahan polinomial dan perkalian skalar yang biasa.
   \end{exparts}
   \textit{Petunjuk.} Gunakan \nearbyexample{PlaneThruOriginSubsp} sebagai
   panduan. Sebagian besar dari kesepuluh syarat hanya memerlukan verifikasi.
   \begin{answer}
      \begin{exparts}
        \partsitem 
          Tiga anggotanya adalah $1+2x$, $2-1x$, dan $x$.
          (Tentu saja, banyak jawaban lain juga mungkin.)

          Verifikasinya sama seperti \nearbyexample{ex:RealVecSpaces}.
          Pertama, kita memeriksa syarat $1$--$5$ dari
          \nearbydefinition{def:VecSpace}, yang berkaitan dengan penjumlahan.
          Untuk ketertutupan terhadap penjumlahan, syarat~(1), perhatikan bahwa
          jika $a+bx,c+dx\in\polyspace_1$, maka
          $(a+bx)+(c+dx)=(a+c)+(b+d)x$ merupakan polinomial linear dengan
          koefisien real, sehingga menjadi anggota $\polyspace_1$.
          Syarat~(2) diperiksa sebagai berikut: jika
          $a+bx,c+dx\in\polyspace_1$, maka
          $(a+bx)+(c+dx)=(a+c)+(b+d)x$, sedangkan dalam urutan sebaliknya
          diperoleh $(c+dx)+(a+bx)=(c+a)+(d+b)x$. Keduanya sama karena
          $a+c=c+a$ dan $b+d=d+b$ bagi bilangan real tersebut.
          Syarat~(3) serupa: misalkan
          $a+bx,c+dx,e+fx\in\polyspace$. Maka
          $((a+bx)+(c+dx))+(e+fx)=(a+c+e)+(b+d+f)x$, sedangkan
          $(a+bx)+((c+dx)+(e+fx))=(a+c+e)+(b+d+f)x$, dan keduanya sama
          (yakni, penjumlahan bilangan real bersifat asosiatif sehingga
          $(a+c)+e=a+(c+e)$ dan $(b+d)+f=b+(d+f)$).
          Untuk syarat~(4), perhatikan bahwa polinomial linear
          $0+0x\in\polyspace_1$ memenuhi
          $(a+bx)+(0+0x)=a+bx$ dan $(0+0x)+(a+bx)=a+bx$.
          Untuk syarat terakhir dalam kelompok ini, yakni syarat~(5),
          perhatikan bahwa bagi setiap $a+bx\in\polyspace_1$, invers aditifnya
          adalah $-a-bx\in\polyspace_1$ karena
          $(a+bx)+(-a-bx)=(-a-bx)+(a+bx)=0+0x$.
         
          Selanjutnya, kita memeriksa syarat~(6)--(10), yang berkaitan dengan
          perkalian skalar.
          Untuk (6), yaitu ketertutupan ruang terhadap perkalian skalar,
          misalkan $r$ bilangan real dan $a+bx$ anggota $\polyspace_1$.
          Maka $r(a+bx)=(ra)+(rb)x$ merupakan anggota $\polyspace_1$ karena
          merupakan polinomial linear dengan koefisien bilangan real.
          Syarat~(7) berlaku karena
          $(r+s)(a+bx)=r(a+bx)+s(a+bx)$ mengikuti sifat distributif perkalian
          bilangan real.
          Syarat~(8) serupa:
          $r((a+bx)+(c+dx))=r((a+c)+(b+d)x)=r(a+c)+r(b+d)x=(ra+rc)+(rb+rd)x
          =r(a+bx)+r(c+dx)$.
          Untuk~(9), kita memperoleh
          $(rs)(a+bx)=(rsa)+(rsb)x=r(sa+sbx)=r(s(a+bx))$.
          Terakhir, syarat~(10) adalah
          $1(a+bx)=(1a)+(1b)x=a+bx$.
        \partsitem 
          Sebut himpunan tersebut $P$.
          Dalam bagian sebelumnya tidak ada pembatasan pada koefisien, tetapi
          di sini kita membatasi perhatian pada polinomial linear dengan
          $a_0-2a_1=0$, yakni suku konstan dikurangi dua kali koefisien suku
          linear bernilai nol.
          Jadi, tiga anggota biasa dari $P$ adalah
          $2+1x$, $6+3x$, dan $-4-2x$.

          Untuk syarat~(1), kita harus menunjukkan bahwa penjumlahan dua
          polinomial linear yang memenuhi pembatasan menghasilkan polinomial
          linear yang juga memenuhi pembatasan. Jika $a+bx,c+dx\in P$, maka
          $(a+bx)+(c+dx)=(a+c)+(b+d)x$ merupakan anggota $P$ karena
          $(a+c)-2(b+d)=(a-2b)+(c-2d)=0+0=0$.
          Syarat~(2) dapat diperiksa sebagai berikut: jika
          $a+bx,c+dx\in\polyspace_1$, maka
          $(a+bx)+(c+dx)=(a+c)+(b+d)x$, sedangkan dalam urutan sebaliknya
          diperoleh $(c+dx)+(a+bx)=(c+a)+(d+b)x$. Keduanya sama karena
          $a+c=c+a$ dan $b+d=d+b$.
          (Jadi, pembatasan tidak memengaruhi syarat ini dan verifikasinya sama
          dengan verifikasi pada bagian pertama latihan ini.)
          Syarat~(3) juga tidak dipengaruhi pembatasan tambahan: jika
          $a+bx,c+dx,e+fx\in\polyspace$, maka
          $((a+bx)+(c+dx))+(e+fx)=(a+c+e)+(b+d+f)x$, sedangkan
          $(a+bx)+((c+dx)+(e+fx))=(a+c+e)+(b+d+f)x$, dan keduanya sama.
          Untuk syarat~(4), perhatikan bahwa polinomial linear $0+0x\in P$
          memenuhi pembatasan karena suku konstannya dikurangi dua kali
          koefisien suku linearnya bernilai nol. Karena itu, verifikasi dari
          bagian pertama berlaku: $0+0x\in\polyspace_1$ memenuhi
          $(a+bx)+(0+0x)=a+bx$ dan $(0+0x)+(a+bx)=a+bx$.
          Untuk memeriksa syarat~(5), perhatikan bahwa bagi setiap $a+bx\in P$,
          invers aditifnya adalah $-a-bx$ karena vektor itu merupakan anggota
          $P$ (dari $a+bx\in P$ kita mengetahui $a-2b=0$, dan mengalikan kedua
          ruas dengan $-1$ menghasilkan $-a+2b=0$). Seperti pada bagian
          pertama, vektor tersebut berperan sebagai invers aditif:
          $(a+bx)+(-a-bx)=(-a-bx)+(a+bx)=0+0x$.
         
          Kita juga harus memeriksa syarat~(6)--(10), yaitu syarat perkalian
          skalar. Untuk (6), ketertutupan ruang terhadap perkalian skalar,
          misalkan $r$ bilangan real dan $a+bx\in P$ (sehingga $a-2b=0$).
          Maka $r(a+bx)=(ra)+(rb)x$ merupakan anggota $P$ karena merupakan
          polinomial linear berkoefisien real yang memenuhi
          $(ra)-2(rb)=r(a-2b)=0$.
          Syarat~(7) berlaku karena alasan yang sama seperti pada bagian pertama
          latihan ini: $(r+s)(a+bx)=r(a+bx)+s(a+bx)$ mengikuti sifat distributif
          perkalian bilangan real.
          Syarat~(8) juga tidak berubah dari bagian pertama:
          $r((a+bx)+(c+dx))=r((a+c)+(b+d)x)=r(a+c)+r(b+d)x=(ra+rc)+(rb+rd)x
          =r(a+bx)+r(c+dx)$.
          Demikian pula~(9):
          $(rs)(a+bx)=(rsa)+(rsb)x=r(sa+sbx)=r(s(a+bx))$.
          Terakhir, demikian pula syarat~(10):
          $1(a+bx)=(1a)+(1b)x=a+bx$.
       \end{exparts}
     \end{answer}
  \item 
    Untuk masing-masing himpunan, daftarkan tiga anggotanya lalu tunjukkan
    bahwa himpunan tersebut merupakan ruang vektor.
    \begin{exparts}
      \partsitem Himpunan matriks \( \nbyn{2} \) dengan entri real di bawah
        operasi matriks yang biasa.
      \partsitem Himpunan matriks \( \nbyn{2} \) dengan entri real dan entri
        $2,1$ bernilai nol, di bawah operasi matriks yang biasa.
    \end{exparts}
    \begin{answer}
      Gunakan \nearbyexample{ex:RealVecSpaces} sebagai panduan.
      (\textit{Catatan.} Karena banyak syarat cukup mudah diperiksa,
      seseorang kadang-kadang merasa pasti ada sesuatu yang terlewat.
      Ingatlah bahwa mudah atau rutin dilakukan tidak sama dengan tidak perlu
      dilakukan.)
      \begin{exparts}
        \partsitem 
          Berikut adalah tiga anggotanya.
          \begin{equation*}
            \begin{mat}
              1  &2  \\
              3  &4
            \end{mat},\,
            \begin{mat}
              -1  &-2  \\
              -3  &-4
            \end{mat},\,
            \begin{mat}
              0  &0  \\
              0  &0
            \end{mat}
          \end{equation*}

          Untuk~(1), jumlah dua matriks real $\nbyn{2}$ merupakan matriks real
          $\nbyn{2}$.
          Untuk~(2), kita meninjau jumlah dua matriks,
          \begin{equation*}
            \begin{mat}
              a  &b  \\
              c  &d
            \end{mat}+
            \begin{mat}
              e  &f  \\
              g  &h
            \end{mat}
            =
            \begin{mat}
              a+e  &b+f  \\
              c+g  &d+h
            \end{mat}
          \end{equation*}
          lalu menerapkan komutativitas penjumlahan bilangan real,
          \begin{equation*}
            =\begin{mat}
              e+a  &f+b  \\
              g+c  &h+d
            \end{mat}
            =
            \begin{mat}
              e  &f  \\
              g  &h
            \end{mat}
            +
            \begin{mat}
              a  &b  \\
              c  &d
            \end{mat}
          \end{equation*}
          untuk memverifikasi bahwa penjumlahan matriks bersifat komutatif.
          Verifikasi syarat~(3), yakni asosiativitas penjumlahan matriks,
          serupa dengan verifikasi sebelumnya:
          \begin{equation*}
            \big(
              \begin{mat}
                a  &b  \\
                c  &d  
              \end{mat}
              +\begin{mat}
                e  &f  \\
                g  &h  
              \end{mat}
            \big)
            +
            \begin{mat}
              i  &j  \\
              k  &l
            \end{mat}
            =
            \begin{mat}
              (a+e)+i  &(b+f)+j  \\
              (c+g)+k  &(d+h)+l
            \end{mat}
          \end{equation*}
          sedangkan
          \begin{equation*}
            \begin{mat}
              a  &b  \\
              c  &d  
            \end{mat}
            +
              \big(
              \begin{mat}
                e  &f  \\
                g  &h  
              \end{mat}
              +
              \begin{mat}
                i  &j  \\
                k  &l
              \end{mat}
            \big)
            =
            \begin{mat}
              a+(e+i)  &b+(f+j)  \\
              c+(g+k)  &d+(h+l)
            \end{mat}
          \end{equation*}
          dan keduanya sama entri demi entri karena penjumlahan bilangan real
          bersifat asosiatif.
          Untuk~(4), unsur nol ruang ini adalah matriks nol $\nbyn{2}$.
          Syarat~(5) berlaku karena bagi setiap matriks $\nbyn{2}$~$A$,
          invers aditifnya adalah matriks yang setiap entrinya merupakan lawan
          entri $A$, yaitu matriks $-1\cdot A$.

          Syarat~(6) berlaku karena kelipatan skalar suatu matriks $\nbyn{2}$
          juga merupakan matriks $\nbyn{2}$.
          Untuk syarat~(7), kita memperoleh
          \begin{multline*}
            (r+s)
            \begin{mat}
              a  &b  \\
              c  &d
            \end{mat}
            =
            \begin{mat}
              (r+s)a  &(r+s)b  \\
              (r+s)c  &(r+s)d
            \end{mat}               \\
            =
            \begin{mat}
              ra+sa  &rb+sb  \\
              rc+sc  &rd+sd
            \end{mat}
            =
            r
            \begin{mat}
              a  &b  \\
              c  &d 
            \end{mat}
            +
            s
            \begin{mat}
              a  &b  \\
              c  &d 
            \end{mat}
          \end{multline*}
          Syarat~(8) diperiksa dengan cara yang sama.
          \begin{multline*}
            r
            \big(
              \begin{mat}
                a  &b  \\
                c  &d
              \end{mat}
              +
              \begin{mat}
                e  &f  \\
                g  &h
              \end{mat}
            \big)
            =
            r
            \begin{mat}
              a+e  &b+f  \\
              c+g  &d+h
            \end{mat}
            =
            \begin{mat}
              ra+re  &rb+rf  \\
              rc+rg  &rd+rh
            \end{mat}              \\
            =
            r
            \begin{mat}
              a  &b  \\
              c  &d 
            \end{mat}
            +
            r
            \begin{mat}
              e  &f  \\
              g  &h 
            \end{mat}
            =
            r
            \big(
            \begin{mat}
              a  &b  \\
              c  &d
            \end{mat}
            +
            \begin{mat}
              e  &f  \\
              g  &h
            \end{mat}
            \big)
          \end{multline*}
          Untuk~(9), kita memperoleh
          \begin{equation*}
            (rs)
            \begin{mat}
              a  &b  \\
              c  &d
            \end{mat}
            =
            \begin{mat}
              rsa  &rsb  \\
              rsc  &rsd
            \end{mat}
            =
            r
            \begin{mat}
              sa  &sb  \\
              sc  &sd
            \end{mat}
            =
            r\big( s
            \begin{mat}
              a  &b  \\
              c  &d
            \end{mat}
            \big)
          \end{equation*}
          Syarat~(10) juga sama mudahnya.
          \begin{equation*}
            1
            \begin{mat}
              a  &b  \\
              c  &d
            \end{mat}
            =
            \begin{mat}
              1\cdot a  &1\cdot b  \\
              1\cdot c  &1\cdot d
            \end{mat}
            =
            \begin{mat}
              a  &b  \\
              c  &d
            \end{mat}
          \end{equation*}
        \partsitem 
          Bagian ini hanya berbeda dari bagian sebelumnya karena kita
          membatasi diri pada himpunan $T$ yang terdiri atas matriks dengan
          entri nol pada baris kedua dan kolom pertama.
          Berikut adalah tiga anggota $T$.
          \begin{equation*}
            \begin{mat}
              1  &2  \\
              0  &4
            \end{mat},\,
            \begin{mat}
              -1  &-2  \\
              0  &-4
            \end{mat},\,
            \begin{mat}
              0  &0  \\
              0  &0
            \end{mat}
          \end{equation*}
          Beberapa verifikasi untuk bagian ini sama dengan verifikasi pada
          bagian pertama latihan ini; di bawah ini kita hanya mengerjakan yang
          berbeda.

          Untuk~(1), jumlah dua matriks real $\nbyn{2}$ dengan entri $2,1$
          bernilai nol juga merupakan matriks real $\nbyn{2}$ dengan entri
          $2,1$ bernilai nol.
          \begin{equation*}
            \begin{mat}
              a  &b  \\
              0  &d
            \end{mat}
            +
            \begin{mat}
              e  &f  \\
              0  &h
            \end{mat}
            =
            \begin{mat}
              a+e  &b+f  \\
              0    &d+h
            \end{mat}
          \end{equation*}
          Verifikasi syarat~(2) yang diberikan pada bagian sebelumnya juga
          berlaku di sini. Hal yang sama berlaku untuk syarat~(3).
          Untuk~(4), perhatikan bahwa matriks nol $\nbyn{2}$ merupakan anggota
          $T$.
          Syarat~(5) berlaku karena invers aditif dari setiap matriks
          $\nbyn{2}$~$A$ adalah matriks $-1\cdot A$; jadi, invers aditif suatu
          matriks yang entri $2,1$-nya nol juga memiliki entri $2,1$ nol.

          Syarat~(6) berlaku karena kelipatan skalar suatu matriks $\nbyn{2}$
          yang entri $2,1$-nya nol juga merupakan matriks $\nbyn{2}$ dengan
          entri $2,1$ nol.
          Verifikasi syarat~(7) sama seperti pada bagian sebelumnya.
          Demikian pula verifikasi syarat~(8), (9), dan~(10).
      \end{exparts}        
    \end{answer}
  \recommended \item 
    Untuk masing-masing himpunan, daftarkan tiga anggotanya lalu tunjukkan
    bahwa himpunan tersebut merupakan ruang vektor.
    \begin{exparts}
      \partsitem Himpunan vektor baris tiga komponen dengan operasi yang biasa.
      \partsitem Himpunan
        \begin{equation*}
          \set{\colvec{x \\ y \\ z \\ w}\in\Re^4\suchthat x+y-z+w=0}
        \end{equation*}
        di bawah operasi yang diwarisi dari $\Re^4$.
    \end{exparts}
    \begin{answer}
      % Most of the conditions are easy to check; use
      % \nearbyexample{ex:RealVecSpaces} as a guide.
      \begin{exparts}
        \partsitem 
          Tiga anggotanya adalah $\rowvec{1  &2  &3}$,
          $\rowvec{2  &1  &3}$, dan $\rowvec{0  &0  &0}$.

          Kita harus memeriksa syarat (1)--(10) dalam
          \nearbydefinition{def:VecSpace}.
          Syarat (1)--(5) berkaitan dengan penjumlahan.
          Untuk syarat~(1), ingat bahwa jumlah dua vektor baris tiga komponen
          \begin{equation*}
            \rowvec{a  &b  &c}
            +\rowvec{d  &e  &f}
            =\rowvec{a+d  &b+e  &c+f}
          \end{equation*}
          juga merupakan vektor baris tiga komponen
          (semua huruf $a,\ldots,f$ menyatakan bilangan real).
          Verifikasi~(2) bersifat rutin,
          \begin{multline*}
            \rowvec{a  &b  &c}
            +\rowvec{d  &e  &f}
            =\rowvec{a+d  &b+e  &c+f}   \\
            =\rowvec{d+a  &e+b  &f+c}
            =\rowvec{d  &e  &f}
            +\rowvec{a  &b  &c}
          \end{multline*}
          (kesamaan kedua berlaku karena ketiga entri merupakan bilangan real
          dan penjumlahan bilangan real bersifat komutatif).
          Verifikasi syarat~(3) serupa.
          \begin{multline*}
            \big(\rowvec{a  &b  &c}
              +\rowvec{d  &e  &f}
            \big)
            +\rowvec{g  &h  &i}
            =\rowvec{(a+d)+g  &(b+e)+h  &(c+f)+i}    \\
            =\rowvec{a+(d+g)  &b+(e+h)  &c+(f+i)}
            =\rowvec{a  &b  &c}
            +\big(\rowvec{d  &e  &f}
               +\rowvec{g  &h  &i}
             \big) 
          \end{multline*}
          Untuk (4), perhatikan bahwa vektor baris tiga komponen
          $\rowvec{0  &0  &0}$ merupakan identitas aditif:
          $\rowvec{a  &b  &c}+\rowvec{0  &0  &0}=\rowvec{a  &b  &c}$.
          Untuk memverifikasi syarat~(5), andaikan diberikan anggota
          $\rowvec{a  &b  &c}$ dari himpunan tersebut. Perhatikan bahwa
          $\rowvec{-a  &-b  &-c}$ juga berada dalam himpunan dan memiliki
          sifat yang diinginkan:
          $\rowvec{a  &b  &c}+\rowvec{-a  &-b  &-c}=\rowvec{0  &0  &0}$.

          Syarat (6)--(10) berkaitan dengan perkalian skalar.
          Untuk memverifikasi~(6), yakni ketertutupan ruang terhadap operasi
          perkalian skalar yang diberikan, perhatikan bahwa
          $r\rowvec{a  &b  &c}=\rowvec{ra  &rb  &rc}$ merupakan vektor baris
          tiga komponen dengan entri real.
          Untuk~(7), kita menghitung
          \begin{multline*}
            (r+s)\rowvec{a  &b  &c}=\rowvec{(r+s)a  &(r+s)b  &(r+s)c}
            =\rowvec{ra+sa  &rb+sb  &rc+sc}                             \\
            =\rowvec{ra  &rb  &rc}+\rowvec{sa  &sb  &sc}
            =r\rowvec{a  &b  &c}+s\rowvec{a  &b  &c}
          \end{multline*}
          Syarat~(8) sangat serupa.
          \begin{multline*} 
            r\big(\rowvec{a  &b  &c}+\rowvec{d  &e  &f}\big)  \\
            \begin{aligned}
            &=r\rowvec{a+d  &b+e  &c+f}
            =\rowvec{r(a+d)  &r(b+e)  &r(c+f)}                     \\
            &=\rowvec{ra+rd  &rb+re  &rc+rf}
            =\rowvec{ra  &rb  &rc}+\rowvec{rd  &re  &rf}          \\
            &=r\rowvec{a  &b  &c}+r\rowvec{d  &e  &f}
            \end{aligned}
          \end{multline*}
          Demikian pula perhitungan untuk syarat~(9).
          \begin{equation*}
            (rs)\rowvec{a  &b  &c}
             =\rowvec{rsa  &rsb  &rsc}
             =r\rowvec{sa  &sb  &sc}
             =r\big(s\rowvec{a  &b  &c}\big)
          \end{equation*}
          Syarat~(10) juga rutin:
          $1\rowvec{a  &b  &c}
             =\rowvec{1\cdot a  &1\cdot b  &1\cdot c}
             =\rowvec{a  &b  &c}$.
        \partsitem 
          Sebut himpunan tersebut $L$.
          Ketertutupan terhadap penjumlahan, syarat~(1), diperiksa dengan
          memastikan bahwa jika kedua suku merupakan anggota~$L$, maka
          jumlahnya,
          \begin{equation*}
            \colvec{a \\ b \\ c \\ d}
            +\colvec{e \\ f \\ g \\ h}
            =
            \colvec{a+e \\ b+f \\ c+g \\ d+h}
          \end{equation*}
          juga merupakan anggota \( L \). Hal ini benar karena jumlah tersebut
          memenuhi kriteria keanggotaan $L$:
          \( (a+e)+(b+f)-(c+g)+(d+h)
          =(a+b-c+d)+(e+f-g+h)=0+0 \).
          Verifikasi syarat~(2), (3), dan~(5) serupa dengan verifikasi pada
          bagian pertama latihan ini.
          Untuk syarat~(4), perhatikan bahwa vektor nol merupakan anggota~$L$
          karena komponen pertama ditambah komponen kedua, dikurangi komponen
          ketiga, lalu ditambah komponen keempat bernilai nol.

          Syarat~(6), ketertutupan terhadap perkalian skalar, juga serupa:
          jika vektor tersebut merupakan anggota~$L$, maka
          \begin{equation*}
            r\colvec{a  \\ b \\ c \\ d}
            =\colvec{ra  \\  rb  \\  rc  \\  rd}
          \end{equation*}
          juga merupakan anggota~$L$ karena
          $ra+rb-rc+rd=r(a+b-c+d)=r\cdot 0=0$.
          Verifikasi syarat~(7), (8), (9), dan~(10) sama seperti pada bagian
          sebelumnya dalam latihan ini.
     \end{exparts}  
     \end{answer}
  \recommended \item \label{exer:NotVectorSpaces}
    Tunjukkan bahwa masing-masing himpunan berikut bukan ruang vektor.
    (\textit{Petunjuk.} Periksa ketertutupan dengan mendaftarkan dua anggota
    setiap himpunan dan mencoba beberapa operasi terhadap keduanya.)
    \begin{exparts}
      \partsitem Di bawah operasi yang diwarisi dari \( \Re^3 \), himpunan
        \begin{equation*}
          \set{\colvec{x \\ y \\ z}\in\Re^3\suchthat x+y+z=1}
        \end{equation*}
      \partsitem Di bawah operasi yang diwarisi dari \( \Re^3 \), himpunan
        \begin{equation*}
          \set{\colvec{x \\ y \\ z}\in\Re^3\suchthat x^2+y^2+z^2=1}
        \end{equation*}
      \partsitem Di bawah operasi matriks yang biasa,
        \begin{equation*}
          \set{\begin{mat}
                 a  &1  \\
                 b  &c
               \end{mat} \suchthat a,b,c\in\Re}
        \end{equation*}
      \partsitem Di bawah operasi polinomial yang biasa,
        \begin{equation*}
          \set{a_0+a_1x+a_2x^2\suchthat a_0,a_1,a_2\in\Re^+}
        \end{equation*}
        dengan $\Re^+$ himpunan bilangan real yang lebih besar dari nol
      \partsitem Di bawah operasi yang diwarisi,
        \begin{equation*}
          \set{\colvec{x \\ y}\in\Re^2\suchthat
               \text{\( x+3y=4 \), \( 2x-y=3 \), dan \( 6x+4y=10 \)} }
        \end{equation*}
    \end{exparts}
    \begin{answer}
      Dalam setiap bagian, himpunannya disebut \( Q \).
      Untuk beberapa bagian, terdapat cara benar lain untuk menunjukkan bahwa
      $Q$ bukan ruang vektor.
      \begin{exparts}
        \partsitem Himpunan ini tidak tertutup terhadap penjumlahan; syarat~(1)
          tidak terpenuhi.
          \begin{equation*}
            \colvec[r]{1 \\ 0 \\ 0},
            \colvec[r]{0 \\ 1 \\ 0}\in Q
            \qquad
            \colvec[r]{1 \\ 1 \\ 0}\not\in Q
          \end{equation*}
        \partsitem Himpunan ini tidak tertutup terhadap penjumlahan.
          \begin{equation*}
            \colvec[r]{1 \\ 0 \\ 0},
            \colvec[r]{0 \\ 1 \\ 0}\in Q
            \qquad
            \colvec[r]{1 \\ 1 \\ 0}\not\in Q
          \end{equation*}
        \partsitem Himpunan ini tidak tertutup terhadap penjumlahan.
          \begin{equation*}
            \begin{mat}[r]
              0  &1  \\
              0  &0
            \end{mat},
            \,
            \begin{mat}[r]
              1  &1  \\
              0  &0
            \end{mat}\in Q
            \qquad
            \begin{mat}[r]
              1  &2  \\
              0  &0
            \end{mat}\not\in Q
          \end{equation*}
        \partsitem Himpunan ini tidak tertutup terhadap perkalian skalar.
          \begin{equation*}
            1+1x+1x^2\in Q
            \qquad
            -1\cdot(1+1x+1x^2)\not\in Q
          \end{equation*}
        \item Himpunan ini kosong sehingga melanggar syarat~(4).
      \end{exparts}  
    \end{answer}
  \item 
    Definisikan operasi penjumlahan dan perkalian skalar yang menjadikan
    bilangan kompleks suatu ruang vektor di atas \( \Re \).
    \begin{answer}
      Operasi yang biasa,
      \( (v_0+v_1i)+(w_0+w_1i)=(v_0+w_0)+(v_1+w_1)i \) dan
      \( r(v_0+v_1i)=(rv_0)+(rv_1)i \), sudah memadai.
      Pemeriksaannya mudah.
    \end{answer}
  \item
    Apakah himpunan bilangan rasional merupakan ruang vektor di atas \( \Re \)
    di bawah operasi penjumlahan dan perkalian skalar yang biasa?
    \begin{answer}
       Tidak. Himpunan tersebut tidak tertutup terhadap perkalian skalar karena,
       misalnya, \( \pi\cdot (1) \) bukan bilangan rasional.
    \end{answer}
  \item 
    Tunjukkan bahwa himpunan kombinasi linear dari variabel \( x,y,z \)
    merupakan ruang vektor di bawah operasi alami penjumlahan dan perkalian
    skalar.
    \begin{answer}
      Operasi alaminya adalah
      \( (v_1x+v_2y+v_3z)+(w_1x+w_2y+w_3z)=(v_1+w_1)x+(v_2+w_2)y+(v_3+w_3)z \)
      dan \( r\cdot(v_1x+v_2y+v_3z)=(rv_1)x+(rv_2)y+(rv_3)z \).
      Pemeriksaan bahwa ini merupakan ruang vektor mudah; gunakan
      \nearbyexample{ex:RealVecSpaces} sebagai panduan.
    \end{answer}
  \item 
    Buktikan bahwa himpunan vektor kolom bertinggi dua dengan entri real di
    bawah operasi berikut bukan ruang vektor.
    \begin{equation*}
      \colvec{x_1 \\ y_1}
      +\colvec{x_2 \\ y_2}
      =\colvec{x_1-x_2 \\ y_1-y_2}
      \qquad
      r\cdot\colvec{x \\ y}
      =\colvec{rx \\ ry}
    \end{equation*}
    \begin{answer}
      Operasi `\( + \)' tidak komutatif (jadi, syarat~(2) tidak terpenuhi);
      mudah memberikan dua anggota himpunan yang membuktikan pernyataan ini.
    \end{answer}
  \item 
    Buktikan atau sangkal bahwa \( \Re^3 \) merupakan ruang vektor di bawah
    operasi-operasi berikut.
    \begin{exparts}
      \partsitem \( 
               \colvec{x_1 \\ y_1 \\ z_1}
               +\colvec{x_2 \\ y_2 \\ z_2}
               =\colvec{0 \\ 0 \\ 0}
               \quad\text{dan}\quad
               r\colvec{x \\ y \\ z}
               =\colvec{rx \\ ry \\ rz} \) 
      \partsitem \( 
               \colvec{x_1 \\ y_1 \\ z_1}
               +\colvec{x_2 \\ y_2 \\ z_2}
               =\colvec{0 \\ 0 \\ 0}   
               \quad\text{dan}\quad
               r\colvec{x \\ y \\ z}
               =\colvec{0 \\ 0 \\ 0} \)
    \end{exparts}
    \begin{answer}
      \begin{exparts}
        \partsitem Himpunan ini bukan ruang vektor.
          \begin{equation*}
            (1+1)\cdot\colvec[r]{1 \\ 0 \\ 0}\neq
            \colvec[r]{1 \\ 0 \\ 0}
            +\colvec[r]{1 \\ 0 \\ 0}
          \end{equation*}
        \partsitem Himpunan ini bukan ruang vektor.
          \begin{equation*}
            1\cdot\colvec[r]{1 \\ 0 \\ 0}\neq\colvec[r]{1 \\ 0 \\ 0}
          \end{equation*}
      \end{exparts}   
    \end{answer}
  \recommended \item
    Untuk masing-masing himpunan, tentukan apakah himpunan tersebut merupakan
    ruang vektor; operasi yang dimaksud adalah operasi alaminya.
    \begin{exparts}
      \partsitem Matriks \( \nbyn{2} \) \definend{diagonal}
        \begin{equation*}
          \set{\begin{mat}
            a  &0  \\
            0  &b
          \end{mat}\suchthat a,b\in\Re}
        \end{equation*}
      \partsitem Himpunan matriks \( \nbyn{2} \) berikut,
        \begin{equation*}
          \set{\begin{mat}
            x    &x+y  \\
            x+y  &y
          \end{mat}\suchthat x,y\in\Re}
        \end{equation*}
      \partsitem Himpunan berikut,
        \begin{equation*}
          \set{\colvec{x \\ y \\ z \\ w}\in\Re^4
               \suchthat x+y+z+w=1}
        \end{equation*}
      \partsitem Himpunan fungsi
        \( \set{\map{f}{\Re}{\Re}\suchthat df/dx+2f=0} \)
      \partsitem Himpunan fungsi
        \( \set{\map{f}{\Re}{\Re}\suchthat df/dx+2f=1} \)
    \end{exparts}
    \begin{answer}
      Untuk setiap jawaban ``ya'', Anda cukup memeriksa beberapa kombinasi
      linear dalam ruang tersebut.
      Untuk setiap jawaban ``tidak'', berikan contoh khusus kegagalan salah
      satu syarat.
      \begin{exparts}
        \partsitem Ya.
          Salah satu kombinasi linear yang dipilih secara bebas adalah
          \begin{equation*}
            2\cdot\begin{mat}
              3  &0 \\
              0  &4
            \end{mat}
            +5\cdot\begin{mat}
              6  &0 \\
              0  &7
            \end{mat}
          \end{equation*}
          Hal kunci yang menjadikan kumpulan tersebut ruang vektor adalah bahwa
          hasilnya merupakan matriks diagonal.
          \begin{equation*}
            =\begin{mat}
              36 &0 \\
              0  &43
            \end{mat}
          \end{equation*}
        \partsitem Ya.
          Berikut adalah suatu kombinasi linear.
          \begin{equation*}
            10\cdot\begin{mat}
              1  &3  \\
              3  &2
            \end{mat}
            +11\cdot\begin{mat}
              3  &7  \\
              7  &4
            \end{mat}
          \end{equation*}
          Kita memeriksa bahwa hasilnya tetap memiliki bentuk dengan jumlah
          entri kiri atas dan kanan bawah sama dengan entri kanan atas serta
          entri kiri bawah.
          \begin{equation*}
            =\begin{mat}
               43  &107  \\
              107  &64
            \end{mat}
          \end{equation*}
        \partsitem Tidak. Himpunan ini tidak tertutup terhadap operasi
          penjumlahan alami. Vektor yang semua komponennya $1/4$ merupakan
          anggota himpunan, tetapi ketika dijumlahkan dengan dirinya sendiri,
          hasilnya, yaitu vektor yang semua komponennya $1/2$, bukan anggota.
        \partsitem Ya.
        \partsitem Tidak. \( f(x)=e^{-2x}+(1/2) \) berada dalam himpunan,
          tetapi \( 2\cdot f \) tidak (jadi, syarat~(6) gagal).
      \end{exparts}  
    \end{answer}
  \recommended \item
    Buktikan atau sangkal bahwa fungsi bernilai real \( f \) dari satu variabel
    real yang memenuhi \( f(7)=0 \) membentuk ruang vektor.
    \begin{answer}
      Himpunan ini merupakan ruang vektor.
      Sebagian besar syarat dalam definisi ruang vektor bersifat rutin; di sini
      kita hanya memeriksa ketertutupan.
      Untuk penjumlahan,
      \( (f_1+f_2)\,(7)=f_1(7)+f_2(7)=0+0=0 \).
      Untuk perkalian skalar,
      \( (r\cdot f)\,(7)=rf(7)=r0=0 \).  
    \end{answer}
  \recommended \item
    Tunjukkan bahwa himpunan \( \Re^+ \) bilangan real positif merupakan ruang
    vektor jika `\( x+y \)' kita tafsirkan sebagai hasil kali \( x \) dan
    \( y \) (sehingga \( 2+3 \) berarti \( 6 \)), dan `\( r\cdot x \)'
    kita tafsirkan sebagai pangkat ke-\( r \) dari \( x \).
    \begin{answer}
      Kita memeriksa \nearbydefinition{def:VecSpace}.

      Pertama, ketertutupan terhadap `\( + \)' berlaku karena hasil kali dua
      bilangan real positif merupakan bilangan real positif.
      Syarat kedua terpenuhi karena perkalian bilangan real bersifat komutatif.
      Demikian pula, karena perkalian bilangan real bersifat asosiatif, syarat
      ketiga terpenuhi.
      Untuk syarat keempat, perhatikan bahwa mengalikan suatu bilangan dengan
      \( 1\in\Re^+ \) tidak mengubah bilangan tersebut.
      Kelima, setiap bilangan real positif memiliki kebalikan yang juga
      merupakan bilangan real positif.

      Syarat keenam, ketertutupan terhadap `\( \cdot \)', berlaku karena setiap
      pangkat bilangan real positif merupakan bilangan real positif.
      Syarat ketujuh hanyalah aturan bahwa \( v^{r+s} \) sama dengan hasil kali
      \( v^r \) dan \( v^s \).
      Syarat kedelapan menyatakan \( (vw)^r=v^rw^r \).
      Syarat kesembilan menyatakan \( (v^r)^s=v^{rs} \).
      Syarat terakhir menyatakan \( v^1=v \).
    \end{answer}
  \item 
      Apakah \( \set{(x,y)\suchthat x,y\in\Re} \) merupakan ruang vektor di
      bawah operasi-operasi berikut?
      \begin{exparts}
        \partsitem \( (x_1,y_1)+(x_2,y_2)=(x_1+x_2,y_1+y_2) \)
        dan \( r\cdot (x,y)=(rx,y) \)
        \partsitem \( (x_1,y_1)+(x_2,y_2)=(x_1+x_2,y_1+y_2) \)
        dan \( r\cdot (x,y)=(rx,0) \)
      \end{exparts}
      \begin{answer}
        \begin{exparts}
           \partsitem Tidak: \( 1\cdot(0,1)+1\cdot(0,1)\neq (1+1)\cdot(0,1) \).
           \partsitem Tidak; perhitungan yang sama seperti pada jawaban
              sebelumnya menunjukkan satu syarat definisi ruang vektor yang
              dilanggar. Contoh pelanggaran lain adalah
              \( 1\cdot (0,1)\neq (0,1) \).
        \end{exparts}  
      \end{answer}
  \item 
    Buktikan atau sangkal bahwa himpunan polinomial berderajat lebih besar dari
    atau sama dengan dua, bersama dengan polinomial nol, merupakan ruang vektor.
    \begin{answer}
      Himpunan ini bukan ruang vektor karena tidak tertutup terhadap
      penjumlahan: \( (x^2)+(1+x-x^2) \) bukan anggota himpunan.
    \end{answer}
  \item 
    Pada tahap ini, istilah ``sama'' baru merupakan intuisi. Meskipun demikian,
    untuk setiap ruang vektor tentukan \( k \) sehingga ruang tersebut ``sama''
    dengan \( \Re^k \).
    \begin{exparts}
      \partsitem Matriks \( \nbym{2}{3} \) di bawah operasi yang biasa
      \partsitem Matriks \( \nbym{n}{m} \) di bawah operasi yang biasa
      \partsitem Himpunan matriks \( \nbyn{2} \) berikut,
        \begin{equation*}
          \set{\begin{mat}
            a  &0  \\
            b  &c
          \end{mat} \suchthat a,b,c\in\Re}
        \end{equation*}
      \partsitem Himpunan matriks \( \nbyn{2} \) berikut,
        \begin{equation*}
          \set{\begin{mat}
            a  &0  \\
            b  &c
          \end{mat} \suchthat a+b+c=0}
        \end{equation*}
    \end{exparts}
    \begin{answer}
      \begin{exparts}
        \partsitem \( 6 \)
        \partsitem \( nm \)
        \partsitem \( 3 \)
        \partsitem Untuk melihat bahwa jawabannya adalah \( 2 \), tulis ulang
          himpunan tersebut sebagai
        \begin{equation*}
          \set{\begin{mat}
            a  &0  \\
            b  &-a-b
          \end{mat} \suchthat a,b\in\Re}
        \end{equation*}
        sehingga terdapat dua parameter.
      \end{exparts}  
     \end{answer}
  \item 
    Dengan menggunakan \( \vec{+} \) untuk menyatakan penjumlahan vektor dan
    \( \,\vec{\cdot}\, \) untuk perkalian skalar, nyatakan ulang definisi
    ruang vektor.
    \begin{answer}
      {
      \def\plus{\mathbin{\vec{+}}}
      \def\tim{\mathbin{\vec{\cdot}}}
        Suatu \definend{ruang vektor}\index{ruang vektor!definisi}
        (di atas \( \Re \)) terdiri atas suatu himpunan \( V \) beserta dua
        operasi `\( \plus \)' dan `\( \tim \)' yang memenuhi syarat-syarat
        berikut.
        Jika \( \vec{v},\vec{w}\in V \), maka
        (1)~\definend{jumlah vektor} keduanya,
          \( \vec{v}\plus\vec{w} \), merupakan anggota \( V \).
        Jika \( \vec{u},\vec{v},\vec{w}\in V \), maka
        (2)~\( \vec{v}\plus\vec{w}=\vec{w}\plus\vec{v} \) dan
        (3)~\( (\vec{v}\plus\vec{w})\plus\vec{u}
                 =\vec{v}\plus(\vec{w}\plus\vec{u}) \).
        (4)~Terdapat suatu \definend{vektor nol}
          \( \zero\in V \) sedemikian sehingga
          \( \vec{v}\plus\zero=\vec{v}\, \) untuk semua \( \vec{v}\in V\).
        (5)~Setiap \( \vec{v}\in V \) memiliki suatu
          \definend{invers aditif} \( \vec{w}\in V \) sedemikian sehingga
          \( \vec{w}\plus\vec{v}=\zero \).
        Jika \( r,s \) merupakan \definend{skalar\/}, yaitu anggota \( \Re \),
        dan \( \vec{v},\vec{w}\in V \), maka
        (6)~setiap \definend{kelipatan skalar}
           \( r\tim\vec{v} \) berada dalam \( V \).
        Jika \( r,s\in\Re \) dan \( \vec{v},\vec{w}\in V \), maka
        (7)~\( (r+s)\tim\vec{v}
           =(r\tim\vec{v})\plus(s\tim\vec{v}) \),
        (8)~\( r\tim (\vec{v}\plus\vec{w})
           =(r\tim\vec{v})\plus(r\tim\vec{w}) \),
        (9)~\( (rs)\tim \vec{v}=r\tim (s\tim\vec{v}) \), dan
        (10)~\( 1\tim \vec{v}=\vec{v} \).
     }
    \end{answer}
  \item 
    Buktikan pernyataan-pernyataan berikut.
    \begin{exparts}
      \partsitem Untuk setiap $\vec{v}\in V$, jika $\vec{w}\in V$ merupakan
        invers aditif dari $\vec{v}$, maka $\vec{v}$ merupakan invers aditif
        dari $\vec{w}$. Jadi, suatu vektor merupakan invers aditif dari setiap
        invers aditif dirinya.
      \partsitem Penjumlahan vektor dapat dicoret dari kiri:~jika
        \( \vec{v},\vec{s},\vec{t}\in V \), maka
        \( \vec{v}+\vec{s}=\vec{v}+\vec{t}\, \) mengakibatkan
        \( \vec{s}=\vec{t} \).
    \end{exparts}
    \begin{answer}
      \begin{exparts}
        \partsitem Misalkan \( V \) suatu ruang vektor, \( \vec{v}\in V \),
          dan andaikan \( \vec{w}\in V \) merupakan invers aditif dari
          $\vec{v}$ sehingga \( \vec{w}+\vec{v}=\zero \).
          Karena penjumlahan komutatif,
          \( \zero=\vec{w}+\vec{v}=\vec{v}+\vec{w} \),
          sehingga \( \vec{v} \) juga merupakan invers aditif dari
          \( \vec{w} \).
        \partsitem Misalkan \( V \) suatu ruang vektor dan
          \( \vec{v},\vec{s},\vec{t}\in V \).
          Invers aditif dari \( \vec{v} \) adalah \( -\vec{v} \), sehingga
          \( \vec{v}+\vec{s}=\vec{v}+\vec{t} \) memberikan
          \( -\vec{v}+\vec{v}+\vec{s}=-\vec{v}+\vec{v}+\vec{t} \),
          yang menyatakan \( \zero+\vec{s}=\zero+\vec{t} \), dan karena itu
          \( \vec{s}=\vec{t} \).
      \end{exparts}  
     \end{answer}
  \item 
    Definisi ruang vektor tidak secara eksplisit menyatakan
    \( \zero+\vec{v}=\vec{v} \) (definisi itu justru menyatakan
    \( \vec{v}+\zero=\vec{v} \)). Tunjukkan bahwa persamaan pertama tetap
    harus berlaku dalam setiap ruang vektor.
    \begin{answer}
      Penjumlahan bersifat komutatif, sehingga dalam setiap ruang vektor dan
      untuk setiap vektor \( \vec{v} \), kita memperoleh
      \( \vec{v}=\vec{v}+\zero=\zero+\vec{v} \).  
    \end{answer}
  \recommended \item
    Buktikan atau sangkal bahwa himpunan semua matriks di bawah operasi yang
    biasa merupakan ruang vektor.
    \begin{answer}
      Himpunan ini bukan ruang vektor karena penjumlahan dua matriks yang
      ukurannya berbeda tidak didefinisikan, sehingga himpunan tersebut gagal
      memenuhi syarat ketertutupan.
    \end{answer}
  \item 
    Dalam ruang vektor, setiap anggota memiliki invers aditif.
    Dapatkah suatu anggota memiliki dua invers aditif atau lebih?
    \begin{answer}
      Setiap anggota ruang vektor memiliki tepat satu invers aditif.

      Misalkan \( V \) suatu ruang vektor dan \( \vec{v}\in V \).
      Jika \( \vec{w}_1,\vec{w}_2\in V \) keduanya merupakan invers aditif
      dari \( \vec{v} \), tinjau \( \vec{w}_1+\vec{v}+\vec{w}_2 \).
      Di satu pihak, ungkapan tersebut sama dengan
      $\vec{w}_1+(\vec{v}+\vec{w}_2)=\vec{w}_1+\zero=\vec{w}_1$.
      Di pihak lain, ungkapan itu sama dengan
      $(\vec{w}_1+\vec{v})+\vec{w}_2=\zero+\vec{w}_2=\vec{w}_2$.
      Oleh karena itu, $\vec{w}_1=\vec{w}_2$.
    \end{answer}
  \item 
    \begin{exparts}
      \partsitem Buktikan bahwa setiap titik, garis, atau bidang yang melalui
         titik asal dalam \( \Re^3 \) merupakan ruang vektor di bawah operasi
         yang diwarisi.
      \partsitem Bagaimana jika himpunan tersebut tidak memuat titik asal?
    \end{exparts}
    \begin{answer}
     \begin{exparts}
      \partsitem Setiap himpunan semacam itu berbentuk
        \( \set{r\cdot\vec{v}+s\cdot\vec{w}\suchthat r,s\in\Re} \)
        dengan salah satu atau kedua vektor \( \vec{v},\vec{w} \) mungkin
        bernilai \( \zero \).
        Di bawah operasi yang diwarisi, ketertutupan terhadap penjumlahan,
        \( (r_1\vec{v}+s_1\vec{w})+(r_2\vec{v}+s_2\vec{w})
           =(r_1+r_2)\vec{v}+(s_1+s_2)\vec{w} \)
        dan perkalian skalar,
        \( c(r\vec{v}+s\vec{w})=(cr)\vec{v}+(cs)\vec{w} \)
        mudah diperiksa. Syarat-syarat lainnya juga rutin.
      \partsitem Himpunan semacam itu tidak dapat menjadi ruang vektor di bawah
        operasi yang diwarisi karena tidak memiliki unsur nol.
     \end{exparts}  
    \end{answer}
  \item 
    Dengan gagasan ruang vektor, kita dapat dengan mudah membuktikan kembali
    bahwa himpunan solusi suatu sistem linear homogen memiliki satu anggota
    atau tak hingga banyak anggota.
    Andaikan \( \vec{v}\in V \) bukan \( \zero \).
    \begin{exparts}
      \partsitem Buktikan bahwa \( r\cdot\vec{v}=\zero \) jika dan hanya jika
        \( r=0 \).
      \partsitem Buktikan bahwa
        \( r_1\cdot\vec{v}=r_2\cdot\vec{v} \) jika dan hanya jika
        \( r_1=r_2 \).
      \partsitem Buktikan bahwa setiap ruang vektor tak trivial bersifat tak
        hingga.
      \partsitem Gunakan fakta bahwa himpunan solusi tak kosong suatu sistem
        linear homogen merupakan ruang vektor untuk menarik kesimpulan.
    \end{exparts}
    \begin{answer}
      Andaikan \( \vec{v}\in V \) bukan \( \zero \).
      \begin{exparts}
        \partsitem Salah satu arah pernyataan jika dan hanya jika sudah jelas:
          jika $r=0$, maka $r\cdot\vec{v}=\zero$.
          Untuk arah sebaliknya, misalkan \( r \) skalar tak nol.
          Jika \( r\vec{v}=\zero \), maka
          \( (1/r)\cdot r\vec{v}=(1/r)\cdot \zero \) menunjukkan bahwa
          $\vec{v}=\zero$, bertentangan dengan pengandaian.
        \partsitem Jika \( r_1,r_2 \) skalar, maka
          \( r_1\vec{v}=r_2\vec{v}\, \) berlaku jika dan hanya jika
          \( (r_1-r_2)\vec{v}=\zero \).
          Menurut bagian sebelumnya, hal ini berarti \( r_1-r_2=0 \).
        \partsitem Ruang tak trivial memiliki suatu vektor
          \( \vec{v}\neq\zero \). Tinjau himpunan
          \( \set{k\cdot\vec{v}\suchthat k\in\Re} \).
          Menurut bagian sebelumnya, himpunan ini tak hingga.
        \partsitem Himpunan solusi bersifat trivial atau tak trivial.
          Dalam kasus kedua, himpunan tersebut tak hingga.
     \end{exparts}  
    \end{answer}
  \item 
    Apakah fungsi bernilai real dari satu variabel real yang terdiferensialkan
    membentuk ruang vektor di bawah operasi alaminya?
    \begin{answer}
      Ya.
      Suatu teorema Kalkulus semester pertama menyatakan bahwa jumlah fungsi
      yang terdiferensialkan juga terdiferensialkan dan
      \( (f+g)^\prime=f^\prime+g^\prime \), serta bahwa kelipatan suatu fungsi
      yang terdiferensialkan juga terdiferensialkan dan
      \( (r\cdot f)^\prime=r\,f^\prime \).
    \end{answer}
  \item 
    Suatu \definend{ruang vektor di atas bilangan kompleks}%
    \index{ruang vektor!skalar kompleks}%
    \index{bilangan kompleks!ruang vektor di atas}
    $\C$ memiliki definisi yang sama dengan ruang vektor di atas bilangan real,
    kecuali bahwa skalar diambil dari \( \C \), bukan \( \Re \).
    Tunjukkan bahwa masing-masing himpunan berikut merupakan ruang vektor di
    atas bilangan kompleks.
    (Ingat cara menjumlahkan dan mengalikan bilangan kompleks:
    \( (a_0+a_1i)+(b_0+b_1i)=(a_0+b_0)+(a_1+b_1)i \) dan
    \( (a_0+a_1i)(b_0+b_1i)=(a_0b_0-a_1b_1)+(a_0b_1+a_1b_0)i \).)
    \begin{exparts}
      \partsitem Himpunan polinomial berderajat~dua dengan koefisien kompleks
      \partsitem Himpunan berikut,
        \begin{equation*}
          \set{\begin{mat}
                 0  &a  \\
                 b  &0
               \end{mat}\suchthat a,b\in\C\text{ dan }
                               a+b=0+0i }
        \end{equation*}
    \end{exparts}
    \begin{answer}
      Pemeriksaannya rutin.
      Perhatikan bahwa `\( 1 \)' berarti \( 1+0i \), dan unsur-unsur nolnya
      adalah berikut.
      \begin{exparts}
        \partsitem \( (0+0i)+(0+0i)x+(0+0i)x^2 \)
        \partsitem \( \begin{mat}
                   0+0i  &0+0i  \\
                   0+0i  &0+0i
                 \end{mat} \)
      \end{exparts}  
    \end{answer}
  \item 
    Sebutkan suatu sifat yang dimiliki semua \( \Re^n \), tetapi tidak
    dicantumkan sebagai syarat ruang vektor.
    \begin{answer}
      Salah satu hal yang tidak terdapat dalam definisi ruang vektor adalah
      ukuran jarak.
    \end{answer}
  \item 
    \begin{exparts}
      \partsitem Buktikan bahwa untuk empat vektor sembarang
        \( \vec{v}_1,\ldots,\vec{v}_4\in V \), kita dapat menempatkan tanda
        kurung dalam jumlahnya dengan cara apa pun tanpa mengubah hasil.
        \begin{multline*}
          ((\vec{v}_1+\vec{v}_2)+\vec{v}_3)+\vec{v}_4
          =(\vec{v}_1+(\vec{v}_2+\vec{v}_3))+\vec{v}_4  
          =(\vec{v}_1+\vec{v}_2)+(\vec{v}_3+\vec{v}_4)  \\
          =\vec{v}_1+((\vec{v}_2+\vec{v}_3)+\vec{v}_4)  
          =\vec{v}_1+(\vec{v}_2+(\vec{v}_3+\vec{v}_4))
        \end{multline*}
        Hal ini memungkinkan kita menulis
        `\( \vec{v}_1+\vec{v}_2+\vec{v}_3+\vec{v}_4 \)'
        tanpa ambiguitas.
      \partsitem Buktikan bahwa dua cara apa pun untuk menempatkan tanda kurung
        dalam jumlah berapa pun banyaknya vektor memberikan hasil yang sama.
        (\textit{Petunjuk.} Gunakan induksi pada banyaknya vektor.)
    \end{exparts}
    \begin{answer}
      \begin{exparts}
        \partsitem Sedikit penataan ulang memberikan hasil yang diinginkan.
          \begin{align*}
            (\vec{v}_1+(\vec{v}_2+\vec{v}_3))+\vec{v}_4
            &=((\vec{v}_1+\vec{v}_2)+\vec{v}_3)+\vec{v}_4  \\
            &=(\vec{v}_1+\vec{v}_2)+(\vec{v}_3+\vec{v}_4)  \\
            &=\vec{v}_1+(\vec{v}_2+(\vec{v}_3+\vec{v}_4))  \\
            &=\vec{v}_1+((\vec{v}_2+\vec{v}_3)+\vec{v}_4)
          \end{align*}
          Setiap kesamaan di atas mengikuti asosiativitas tiga vektor yang
          diberikan sebagai salah satu syarat dalam definisi ruang vektor.
          Sebagai contoh, tanda `$=$' kedua menerapkan aturan
          $(\vec{w}_1+\vec{w}_2)+\vec{w}_3=\vec{w}_1+(\vec{w}_2+\vec{w}_3)$
          dengan mengambil $\vec{w}_1=\vec{v}_1+\vec{v}_2$,
          $\vec{w}_2=\vec{v}_3$, dan $\vec{w}_3=\vec{v}_4$.
        \partsitem Kasus dasar induksi adalah kasus tiga vektor.
          Persamaan
          \( \vec{v}_1+(\vec{v}_2+\vec{v}_3)
          =(\vec{v}_1+\vec{v}_2)+\vec{v}_3 \) merupakan salah satu syarat
          dalam definisi ruang vektor.

          Untuk langkah induksi, andaikan setiap dua jumlah tiga vektor,
          setiap dua jumlah empat vektor, \ldots, dan setiap dua jumlah
          $k$ vektor bernilai sama, apa pun penempatan tanda kurungnya.
          Kita akan menunjukkan bahwa setiap jumlah \( k+1 \) vektor sama
          dengan jumlah berikut,
          \( ((\cdots((\vec{v}_1+\vec{v}_2)+\vec{v}_3)+\cdots)+\vec{v}_k)
              +\vec{v}_{k+1} \).

          Setiap jumlah bertanda kurung memiliki `\( + \)' terluar.
          Andaikan tanda tersebut berada di antara \( \vec{v}_m \) dan
          \( \vec{v}_{m+1} \), sehingga jumlahnya berbentuk berikut.
          \begin{equation*}
            (\cdots\,\vec{v}_1\cdots\vec{v}_m\,\cdots)
           +(\cdots\,\vec{v}_{m+1}\cdots\vec{v}_{k+1}\,\cdots)
          \end{equation*}
          Paruh kedua melibatkan kurang dari $k+1$ penjumlahan, sehingga
          menurut hipotesis induksi kita dapat menempatkan ulang tanda
          kurungnya agar dibaca dari kiri ke kanan, dari bagian dalam ke luar.
          Secara khusus, `$+$' terluarnya dapat ditempatkan tepat sebelum
          $\vec{v}_{k+1}$.
          \begin{equation*}
            =(\cdots\,\vec{v}_1\,\cdots\,\vec{v}_m\,\cdots)
             +((\cdots(\vec{v}_{m+1}+\vec{v}_{m+2})+\cdots+\vec{v}_{k})
                 +\vec{v}_{k+1})
          \end{equation*}
          Terapkan asosiativitas jumlah tiga objek,
          \begin{equation*}
            =((\,\cdots\, \vec{v}_1\,\cdots\,\vec{v}_m\,\cdots\,)
             +(\,\cdots\,(\vec{v}_{m+1}+\vec{v}_{m+2})+\cdots\,\vec{v}_k))
             +\vec{v}_{k+1}
          \end{equation*}
          lalu selesaikan dengan menerapkan hipotesis induksi di dalam tanda
          kurung terluar tersebut.
      \end{exparts}  
    \end{answer}
  \item \nearbyexample{ex:ColsIntEntNotVS} memberikan subhimpunan $\Re^2$
   yang bukan ruang vektor di bawah operasi
   yang wajar karena, walaupun tertutup terhadap penjumlahan, himpunan itu
   tidak tertutup terhadap perkalian skalar.
   Tinjau himpunan vektor di bidang yang kedua komponennya bertanda sama atau
   bernilai~$0$.
   Tunjukkan bahwa himpunan ini tertutup terhadap perkalian skalar, tetapi tidak
   terhadap penjumlahan.
   \begin{answer}
     Misalkan $\vec{v}$ anggota $\Re^2$ dengan komponen $v_1$ dan $v_2$.
     Syarat bahwa kedua komponen bertanda sama atau bernilai~$0$ dapat kita
     singkat sebagai $v_1v_2\geq 0$.

     Untuk menunjukkan bahwa himpunan tertutup terhadap perkalian skalar,
     perhatikan bahwa komponen-komponen $r\vec{v}$ memenuhi
     $(rv_1)(rv_2)=r^2(v_1v_2)$, dan $r^2\geq 0$, sehingga
     $r^2v_1v_2\geq 0$.

     Untuk menunjukkan bahwa himpunan tidak tertutup terhadap penjumlahan,
     kita cukup memberikan satu contoh.
     Vektor dengan komponen $-1$ dan~$0$, ketika dijumlahkan dengan vektor
     berkomponen $0$ dan~$1$, menghasilkan vektor dengan komponen berbeda
     tanda, yaitu $-1$ dan~$1$.
   \end{answer}
  % \item 
  %   For any vector space, a subset that is itself a vector space
  %   under the inherited operations
  %   (e.g., a plane through the origin inside of \( \Re^3 \))
  %   is a \definend{subspace}.
  %   \begin{exparts}
  %     \partsitem Show that \( \set{a_0+a_1x+a_2x^2\suchthat a_0+a_1+a_2=0} \)
  %       is a subspace of the vector space of degree~two polynomials.
  %     \partsitem Show that this is a subspace of the \( \nbyn{2} \) matrices.
  %       \begin{equation*}
  %         \set{\begin{mat}
  %                a  &b  \\
  %                c  &0
  %              \end{mat} \suchthat a+b=0}
  %       \end{equation*}
  %     \partsitem Show that a nonempty subset \( S \) of a real vector
  %       space is a
  %       subspace if and only if it is closed under linear combinations of
  %       pairs of vectors: whenever \( c_1,c_2\in\Re \) and
  %       \( \vec{s}_1,\vec{s}_2\in S \) then the combination
  %       \( c_1\vec{v}_1+c_2\vec{v}_2 \) is in \( S \).
  %   \end{exparts}
  %   \begin{answer}
  %     \begin{exparts}
  %       \partsitem We outline the check of the conditions from
  %         \nearbydefinition{def:VecSpace}.

  %         Additive closure holds because if \( a_0+a_1+a_2=0 \)
  %         and \( b_0+b_1+b_2=0 \) then
  %         \begin{equation*}
  %           (a_0+a_1x+a_2x^2)+(b_0+b_1x+b_2x^2)=
  %           (a_0+b_0)+(a_1+b_1)x+(a_2+b_2)x^2
  %         \end{equation*}
  %         is in the set since
  %         \( (a_0+b_0)+(a_1+b_1)+(a_2+b_2)=(a_0+a_1+a_2)+(b_0+b_1+b_2) \)
  %         is zero.
  %         The second through fifth conditions are easy.

  %         Closure under scalar multiplication holds because
  %         if \( a_0+a_1+a_2=0 \) then
  %         \begin{equation*}
  %           r\cdot(a_0+a_1x+a_2x^2)=
  %           (ra_0)+(ra_1)x+(ra_2)x^2
  %         \end{equation*}
  %         is in the set as \( ra_0+ra_1+ra_2=r(a_0+a_1+a_2) \) is zero.
  %         The remaining conditions here are also easy.
  %       \partsitem This is similar to the prior answer.
  %       \partsitem Call the vector space \( V \).
  %         We have two implications: left to right, if \( S \) is a subspace
  %         then it is closed under linear combinations of pairs of vectors and,
  %         right to left, if a nonempty subset is closed under linear
  %         combinations of pairs of vectors then it is a subspace.
  %         The left to right implication is easy; we here sketch the
  %         other one by
  %         assuming \( S \) is nonempty and closed, and checking the conditions
  %         of \nearbydefinition{def:VecSpace}.

  %         First, to show closure under addition, if
  %         \( \vec{s}_1,\vec{s}_2\in S \) then \( \vec{s}_1+\vec{s}_2\in S \)
  %         as \( \vec{s}_1+\vec{s}_2=1\cdot\vec{s}_1+1\cdot\vec{s}_2 \).
  %         Second, for any \( \vec{s}_1,\vec{s}_2\in S \), because addition
  %         is inherited from \( V \), the sum \( \vec{s}_1+\vec{s}_2 \)
  %         in \( S \) equals the sum \( \vec{s}_1+\vec{s}_2 \)
  %         in \( V \) and that equals the sum \( \vec{s}_2+\vec{s}_1 \) in
  %         \( V \) and that in turn equals the sum \( \vec{s}_2+\vec{s}_1 \) in
  %         \( S \).
  %         The argument for the third condition is similar to that for the
  %         second.
  %         For the fourth, suppose that 
  %         \( \vec{s} \) is in the nonempty set \( S \)
  %         and note that \( 0\cdot\vec{s}=\zero\in S \); showing that 
  %         the \( \zero \) of
  %         \( V \) acts under the inherited operations as the additive identity
  %         of \( S \) is easy.
  %         The fifth condition is satisfied because for any \( \vec{s}\in S \)
  %         closure under linear combinations shows that the vector
  %         \( 0\cdot\zero+(-1)\cdot\vec{s} \) is in \( S \); showing 
  %         that it is the
  %         additive inverse of \( \vec{s} \) under the inherited operations is
  %         routine.

  %         The proofs for the remaining conditions are similar.
  %     \end{exparts}  
  %   \end{answer}
\end{exercises}





















 
\subsection{Subruang dan Himpunan Perentang}
\index{subruang|(}
Dalam \nearbyexample{PlaneThruOriginSubsp}, kita melihat suatu ruang vektor
yang merupakan subhimpunan $\Re^2$, yakni garis yang melalui titik asal.
Di sana, ruang vektor $\Re^2$ memuat ruang vektor lain, yaitu garis tersebut.

\begin{definition} \label{df:Subspace}
%<*df:Subspace>
Untuk setiap ruang vektor, suatu \definend{subruang}%
\index{ruang vektor!subruang}\index{subruang!definisi}
adalah subhimpunan yang dengan sendirinya merupakan ruang vektor di bawah
operasi yang diwarisi.
%</df:Subspace>
\end{definition}

\begin{example}  \label{ex:PlaneSubspRThree}
Bidang yang melalui titik asal ini,
\begin{equation*}
  P=\set{\colvec{x \\ y \\ z}\suchthat x+y+z=0}
\end{equation*}
merupakan subruang dari \( \Re^3 \).
Sebagaimana disyaratkan definisi, operasi pada bidang tersebut diwarisi dari
ruang yang lebih besar; jadi, vektor-vektor dijumlahkan dalam $P$ sebagaimana
dalam $\Re^3$,
\begin{equation*}
   \colvec{x_1 \\ y_1 \\ z_1}+\colvec{x_2 \\ y_2 \\ z_2}
   =\colvec{x_1+x_2 \\ y_1+y_2 \\ z_1+z_2}
\end{equation*}
dan perkalian skalarnya juga sama seperti dalam $\Re^3$.
Untuk menunjukkan bahwa $P$ merupakan subruang, kita hanya perlu memperhatikan
bahwa $P$ adalah subhimpunan lalu memverifikasi bahwa $P$ merupakan ruang.
Kita tidak akan memeriksa kesepuluh syarat, melainkan hanya dua syarat
ketertutupan.
Untuk ketertutupan terhadap penjumlahan, perhatikan bahwa jika kedua suku
memenuhi $x_1+y_1+z_1=0$ dan $x_2+y_2+z_2=0$, maka jumlahnya memenuhi
$(x_1+x_2)+(y_1+y_2)+(z_1+z_2)=(x_1+y_1+z_1)+(x_2+y_2+z_2)=0$.
Untuk ketertutupan terhadap perkalian skalar, jika $x+y+z=0$, maka kelipatan
skalarnya memenuhi
$rx+ry+rz=r(x+y+z)=0$.
\end{example}

\begin{example}   \label{ex:SubspacesRTwo}
Sumbu \( x \) dalam \( \Re^2 \) merupakan subruang dengan operasi penjumlahan
dan perkalian skalar yang diwarisi.
\begin{equation*}
  \colvec{x_1 \\ 0}
    +
  \colvec{x_2 \\ 0}
    =
  \colvec{x_1+x_2 \\ 0}
  \qquad
  r\cdot\colvec{x \\ 0}
  =\colvec{rx \\ 0}
\end{equation*}
Seperti dalam contoh sebelumnya, untuk memverifikasi langsung dari definisi
bahwa himpunan ini merupakan subruang, kita cukup mencatat bahwa himpunan ini
merupakan subhimpunan lalu memeriksa bahwa syarat-syarat definisi ruang vektor
terpenuhi.
Sebagai contoh, kedua syarat ketertutupan terpenuhi: menjumlahkan dua vektor
dengan komponen kedua nol menghasilkan vektor dengan komponen kedua nol, dan
mengalikan vektor semacam itu dengan skalar juga menghasilkan vektor dengan
komponen kedua nol.
\end{example}

\begin{example}
Subruang lain dari $\Re^2$ adalah subruang trivialnya.
\begin{equation*}
  \set{\colvec[r]{0 \\ 0}}
\end{equation*}
\end{example}

%<*ProperSubspace>
Setiap ruang vektor memiliki subruang trivial
\( \set{\zero\,} \).\index{subruang!trivial}%
\index{subruang trivial}
Pada ekstrem sebaliknya, setiap ruang vektor memuat dirinya sendiri sebagai
subruang. Subruang yang bukan seluruh ruang disebut
subruang \definend{sejati}.\index{subruang!sejati}%
\index{subruang sejati}
%</ProperSubspace>

\begin{example}  \label{ex:LinSubspPolyThree}
Ruang vektor yang bukan $\Re^n$ juga memiliki subruang.
Ruang polinomial kubik,
\( \set{a+bx+cx^2+dx^3\suchthat a,b,c,d\in\Re} \) 
memiliki subruang yang terdiri atas semua polinomial linear,
\( \set{m+nx\suchthat m,n\in\Re} \).
\end{example}

\begin{example}
Contoh lain subruang yang bukan subhimpunan suatu $\Re^n$ muncul setelah
definisi ruang vektor.
Ruang semua fungsi bernilai real dari satu variabel real dalam
\nearbyexample{ex:RealValuedFcns},
\( \set{f\suchthat \map{f}{\Re}{\Re} } \), memiliki subruang dalam
\nearbyexample{ex:SpaceSatisfyingDiffQ}, yaitu fungsi-fungsi yang memenuhi
syarat $(d^2\,f/dx^2)+f=0$.
\end{example}

\begin{example} \label{ex:OperNotInherit}
Definisi mensyaratkan bahwa operasi penjumlahan dan perkalian skalar harus
merupakan operasi yang diwarisi dari ruang yang lebih besar.
Himpunan \( S=\set{1} \) merupakan subhimpunan \( \Re^1 \).
Di bawah operasi $1+1=1$ dan $r\cdot 1=1$, himpunan $S$ merupakan ruang
vektor, khususnya ruang trivial.
Namun, $S$ bukan subruang dari \( \Re^1 \) karena operasi-operasi tersebut
bukan operasi yang diwarisi; tentu saja, dalam \( \Re^1 \) berlaku
\( 1+1=2 \).
\end{example}

\begin{example}  \label{cex:RPlusNotSubSp}
Karena subruang sendiri merupakan ruang vektor, subruang harus memenuhi syarat
ketertutupan.
Himpunan \( \Re^+ \) bukan subruang dari ruang vektor \( \Re^1 \) karena di
bawah operasi yang diwarisi, himpunan tersebut tidak tertutup terhadap
perkalian skalar: jika \( \vec{v}=1 \), maka
\( -1\cdot\vec{v}\not\in\Re^+ \).
\end{example}

Hasil berikut menyatakan bahwa \nearbyexample{cex:RPlusNotSubSp} merupakan
contoh prototipe. Jika suatu subhimpunan tak kosong menggunakan operasi yang
diwarisi, satu-satunya alasan subhimpunan itu gagal menjadi subruang adalah
karena tidak tertutup.

\begin{lemma}     \label{th:SubspIffClosed} \index{subruang!tertutup}
%<*lm:SubspIffClosed>
Untuk subhimpunan tak kosong \( S \) dari suatu ruang vektor, di bawah operasi
yang diwarisi, pernyataan-pernyataan berikut ekuivalen.
\appendrefs{ekuivalensi pernyataan}   %\spacefactor=1000
\begin{tfae}
  \item \( S \) merupakan subruang dari ruang vektor tersebut
  \item \( S \) tertutup terhadap kombinasi linear pasangan vektor: untuk
    sebarang vektor \( \vec{s}_1,\vec{s}_2\in S \) dan skalar \( r_1,r_2 \),
    vektor \( r_1\vec{s}_1+r_2\vec{s}_2 \) berada dalam \( S \)
  \item \( S \) tertutup terhadap kombinasi linear berapa pun banyaknya
    vektor: untuk sebarang vektor \( \vec{s}_1,\ldots,\vec{s}_n\in S \) dan
    skalar \( r_1, \ldots,r_n \), vektor
    \( r_1\vec{s}_1+\cdots+r_n\vec{s}_n \) merupakan anggota \( S \).
\end{tfae}
%</lm:SubspIffClosed>
\end{lemma}
\noindent Singkatnya, suatu subhimpunan merupakan subruang jika dan hanya jika
subhimpunan tersebut tertutup terhadap kombinasi linear.

\begin{proof}
%<*pf:SubspIffClosed0>
`Pernyataan-pernyataan berikut ekuivalen' berarti bahwa setiap pasangan
pernyataan saling ekuivalen.
\begin{equation*}
  (1)\!\iff\!(2)
  \qquad
  (2)\!\iff\!(3)
  \qquad
  (3)\!\iff\!(1)
\end{equation*}
Kita akan membuktikan ekuivalensi dengan menunjukkan
\( (1)\implies (3)\implies (2)\implies (1)\).
Strategi ini disarankan oleh pengamatan bahwa implikasi
\( (1)\implies (3) \) dan \( (3)\implies (2) \) mudah, sehingga kita hanya
perlu membuktikan \( (2)\implies (1) \).
%</pf:SubspIffClosed0>

%<*pf:SubspIffClosed1>
Andaikan \( S \) subhimpunan tak kosong dari suatu ruang vektor $V$ yang
tertutup terhadap kombinasi pasangan vektor.
Kita akan menunjukkan bahwa $S$ merupakan ruang vektor dengan memeriksa
syarat-syaratnya.

Definisi ruang vektor memiliki lima syarat mengenai penjumlahan.
Pertama, untuk ketertutupan terhadap penjumlahan, jika
\( \vec{s}_1,\vec{s}_2\in S \), maka \( \vec{s}_1+\vec{s}_2\in S \) karena
jumlah tersebut merupakan kombinasi sepasang vektor dan kita mengandaikan
bahwa~\( S \) tertutup terhadap kombinasi semacam itu.
Kedua, untuk sebarang \( \vec{s}_1,\vec{s}_2\in S \), karena penjumlahan
diwarisi dari \( V \), jumlah \( \vec{s}_1+\vec{s}_2 \) dalam \( S \) sama
dengan jumlah \( \vec{s}_1+\vec{s}_2 \) dalam \( V \). Jumlah itu sama dengan
\( \vec{s}_2+\vec{s}_1 \) dalam \( V \) (karena $V$ merupakan ruang vektor,
penjumlahannya komutatif), yang pada gilirannya sama dengan jumlah
\( \vec{s}_2+\vec{s}_1 \) dalam \( S \).
Argumen untuk syarat ketiga serupa dengan argumen untuk syarat kedua.
Untuk syarat keempat, tinjau vektor nol dari \( V \) dan perhatikan bahwa
ketertutupan $S$ terhadap kombinasi linear pasangan vektor memberikan
\( 0\cdot\vec{s}+0\cdot\vec{s}=\zero \) sebagai anggota \( S\)
(dengan \( \vec{s} \) anggota sembarang dari himpunan tak kosong \( S \)).
Mudah diperiksa bahwa \( \zero \) berperan sebagai identitas aditif \( S \)
di bawah operasi yang diwarisi.
Syarat kelima terpenuhi karena, untuk setiap \( \vec{s}\in S \), ketertutupan
terhadap kombinasi linear pasangan vektor menunjukkan bahwa
\( 0\cdot\zero+(-1)\cdot\vec{s} \) merupakan anggota~\( S \), dan jelas
merupakan invers aditif dari \( \vec{s} \) di bawah operasi yang diwarisi.
%</pf:SubspIffClosed1>
Verifikasi syarat-syarat perkalian skalar serupa; lihat
\nearbyexercise{exer:SubspIffClosed}.
\end{proof}

Biasanya kita akan memverifikasi bahwa suatu subhimpunan merupakan subruang
dengan memeriksa bahwa subhimpunan tersebut memenuhi pernyataan~(2).

\begin{remark}
Pada awal bab ini, kita memperkenalkan ruang vektor sebagai kumpulan tempat
kombinasi linear ``bermakna.''
% The vector space definition has ten conditions but eight of them, the
% conditions not about closure, simply ensure that referring to the
% operations by the names `addition' and `scalar multiplication' is sensible.
% For instance, calling an operation `$+$' would not
% be odd unless it is commutative, as in condition~(2). 
% The proof above checks that the subspace satisfies these 
% eight non-closure conditions because they hold in the
% surrounding vector space, provided that the nonempty set $S$ satisfies
% \nearbylemma{th:SubspIffClosed}'s statement~(2) 
% (e.g., commutativity of addition in $S$ follows right from
% commutativity of addition in $V$).
% There is a second way in which we can regard vector spaces as structures in 
% which linear combinations are ``sensible.''
% It has to do with the closure conditions.
Pernyataan (1)--(3) dalam \nearbylemma{th:SubspIffClosed} menyatakan bahwa
kita selalu dapat memaknai ungkapan seperti
$r_1\vec{s}_1+r_2\vec{s}_2$, dalam arti bahwa vektor yang dideskripsikan
berada dalam himpunan~$S$.

Sebagai pembanding, tinjau himpunan $T$ yang terdiri atas vektor bertinggi dua
dengan jumlah entri lebih besar dari atau sama dengan nol.
Di sini kita tidak dapat begitu saja menulis kombinasi linear sembarang seperti
$2\vec{t}_1-3\vec{t}_2$ dan yakin bahwa hasilnya merupakan anggota $T$.
\end{remark}

\nearbylemma{th:SubspIffClosed} menunjukkan bahwa cara yang baik untuk
memikirkan ruang vektor adalah sebagai kumpulan kombinasi linear tanpa
pembatasan. Dua contoh berikut mengambil beberapa ruang dan menyusun ulang
deskripsinya ke dalam bentuk tersebut.

\begin{example}
Kita dapat menunjukkan bahwa bidang melalui titik asal berikut, yang merupakan
subhimpunan $\Re^3$,
\begin{equation*}
  S=\set{\colvec{x \\ y \\ z}\suchthat x-2y+z=0}
\end{equation*}
merupakan subruang di bawah operasi biasa penjumlahan dan perkalian skalar
vektor kolom dengan memeriksa bahwa himpunan tersebut tak kosong dan tertutup
terhadap kombinasi linear dua vektor.
Namun, ada cara lain.
Pandang $x-2y+z=0$ sebagai sistem linear satu persamaan dan parametrisasi
sistem itu dengan menyatakan variabel utama dalam bentuk variabel-variabel
bebas, $x=2y-z$.
\begin{equation*}
     S
     =\set{\colvec{2y-z \\ y \\ z}\suchthat y,z\in\Re}
     =\set{y\colvec[r]{2 \\ 1 \\ 0}+
            z\colvec[r]{-1 \\ 0 \\ 1}\suchthat y,z\in\Re}
    \tag{$*$}
\end{equation*}
Sekarang, untuk menunjukkan bahwa himpunan ini merupakan subruang, tinjau
$r_1\vec{s}_1+r_2\vec{s}_2$.
Setiap $\vec{s}_i$ merupakan kombinasi linear dari dua vektor dalam ($*$),
sehingga ungkapan tersebut merupakan kombinasi linear dari kombinasi linear.
\begin{equation*}
  r_1\cdot(y_1\colvec[r]{2 \\ 1 \\ 0}+
            z_1\colvec[r]{-1 \\ 0 \\ 1})
  +
  r_2\cdot(y_2\colvec[r]{2 \\ 1 \\ 0}+
            z_2\colvec[r]{-1 \\ 0 \\ 1})
\end{equation*}
Lema Kombinasi Linear, Lema~I.\ref{lm:LinearCombinationLemma}, menunjukkan
bahwa hasil keseluruhannya merupakan kombinasi linear dari kedua vektor itu,
sehingga pernyataan~(2) dalam \nearbylemma{th:SubspIffClosed} terpenuhi.
\end{example}

\begin{example} \label{ex:ParamSubspace}
Himpunan berikut merupakan subruang dari ruang matriks \( \nbyn{2} \),
$\matspace_{\nbyn{2}}$.
\begin{equation*}
  L=\set{\begin{mat}
         a  &0  \\
         b  &c
       \end{mat}
       \suchthat a+b+c=0}
\end{equation*}
Untuk melakukan parametrisasi, nyatakan syarat tersebut sebagai $a=-b-c$.
\begin{equation*}
  L
  =\set{\begin{mat}
         -b-c  &0  \\
         b     &c
       \end{mat}
       \suchthat b,c\in\Re}
  =\set{b\begin{mat}[r]
         -1    &0  \\
         1     &0
       \end{mat}
       +c\begin{mat}[r]
         -1    &0  \\
         0     &1
       \end{mat}
       \suchthat b,c\in\Re}
\end{equation*}
Seperti di atas, kita telah mendeskripsikan subruang sebagai kumpulan kombinasi
linear tanpa pembatasan.
Untuk menunjukkan bahwa himpunan tersebut merupakan subruang, perhatikan bahwa
kombinasi linear vektor-vektor dari $L$ merupakan kombinasi linear dari
kombinasi-kombinasi linear, sehingga pernyataan~(2) benar.
\end{example}

% Parametrization is easy, but important.
% We shall use it often.

\begin{definition} \label{df:Span}
%<*df:Span>
\definend{Rentang}\index{rentang}\index{himpunan!rentang}\index{ketertutupan}
(atau \definend{penutupan linear}) dari suatu subhimpunan tak kosong \( S \)
dalam ruang vektor adalah himpunan semua kombinasi linear vektor-vektor dari
\( S \).
\begin{equation*}
  \spanof{S} =\{ c_1\vec{s}_1+\cdots+c_n\vec{s}_n
            \suchthat c_1,\ldots, c_n\in\Re
            \text{ dan } \vec{s}_1,\ldots,\vec{s}_n\in S \}
\end{equation*}
Rentang subhimpunan kosong suatu ruang vektor adalah subruang trivialnya.
%</df:Span>
\end{definition}

%<*NotationForSpan>
\noindent Tidak ada notasi rentang yang sepenuhnya baku.
Kurung siku yang digunakan di sini lazim, tetapi demikian pula
`$\mbox{span}(S)$' dan `$\mbox{sp}(S)$'.
%</NotationForSpan>

\begin{remark}
Dalam Bab Satu, setelah menunjukkan bahwa kita dapat menulis himpunan solusi
suatu sistem linear homogen sebagai
$\set{c_1\vec{\beta}_1+\cdots+c_k\vec{\beta}_k\suchthat
  c_1,\ldots,c_k\in\Re}$,
kita mendeskripsikannya sebagai himpunan yang `dihasilkan' oleh
vektor-vektor $\smash{\vec{\beta}}$.
Sekarang kita menyebutnya rentang dari
$\set{\vec{\beta}_1,\ldots,\vec{\beta}_k}$.

Ingat pula dari pembuktian tersebut bahwa rentang himpunan kosong didefinisikan
sebagai \( \set{\zero} \), mengikuti konvensi bahwa kombinasi linear trivial,
yakni kombinasi tanpa vektor, berjumlah~\( \zero \).
Selain itu, mendefinisikan rentang himpunan kosong sebagai subruang trivial
memudahkan kita karena hasil seperti yang berikut tidak memerlukan pengecualian
untuk himpunan kosong.
\end{remark}

\begin{lemma}   \label{le:SpanIsASubsp}
%<*lm:SpanIsASubsp>
Dalam suatu ruang vektor, rentang setiap subhimpunan merupakan subruang.
%</lm:SpanIsASubsp>
\end{lemma}

\begin{proof}
%<*pf:SpanIsASubsp>
Jika subhimpunan \( S \) kosong, berdasarkan definisi rentangnya merupakan
subruang trivial.
Jika \( S\) tak kosong, menurut \nearbylemma{th:SubspIffClosed} kita hanya
perlu memeriksa bahwa rentang \( \spanof{S} \) tertutup terhadap kombinasi
linear pasangan anggotanya.
Untuk sepasang vektor dari rentang tersebut,
\( \vec{v}=c_1\vec{s}_1+\cdots+c_n\vec{s}_n \) dan
\( \vec{w}=c_{n+1}\vec{s}_{n+1}+\cdots+c_m\vec{s}_m \),
suatu kombinasi linear
\begin{multline*}
  p\cdot(c_1\vec{s}_1+\cdots+c_n\vec{s}_n)+
       r\cdot(c_{n+1}\vec{s}_{n+1}+\cdots+c_m\vec{s}_m)  \\
  =
  pc_1\vec{s}_1+\cdots+pc_n\vec{s}_n
    +rc_{n+1}\vec{s}_{n+1}+\cdots+rc_m\vec{s}_m
\end{multline*}
merupakan kombinasi linear anggota-anggota~\( S \), sehingga menjadi anggota
\( \spanof{S} \)
(mungkin beberapa $\vec{s}_i$ dari $\vec{v}$ sama dengan beberapa
$\vec{s}_j$ dari $\vec{w}$, tetapi hal itu tidak menjadi masalah).
%</pf:SpanIsASubsp>
\end{proof}

Kebalikan lema tersebut juga berlaku: setiap subruang merupakan rentang suatu
himpunan karena suatu subruang jelas merupakan rentang dirinya sendiri, yaitu
himpunan semua anggotanya.
Jadi, suatu subhimpunan ruang vektor merupakan subruang jika dan hanya jika
subhimpunan tersebut merupakan suatu rentang.
Hal ini sesuai dengan intuisi bahwa ruang vektor sebaiknya dipandang sebagai
kumpulan tempat kombinasi linear dapat dibentuk dengan masuk akal.

Secara bersama-sama, \nearbylemma{th:SubspIffClosed} dan
\nearbylemma{le:SpanIsASubsp} menunjukkan bahwa rentang subhimpunan $S$ dari
suatu ruang vektor adalah subruang terkecil yang memuat semua anggota $S$.

\begin{example}   \label{ex:SpanSingVec}
Dalam setiap ruang vektor \( V \), bagi setiap vektor \( \vec{v}\in V \),
himpunan \( \set{r\cdot\vec{v} \suchthat r\in\Re} \) merupakan subruang
dari \( V \).
Sebagai contoh, untuk setiap vektor \( \vec{v}\in\Re^3 \), garis yang melalui
titik asal dan memuat vektor tersebut,
\( \set{k\vec{v}\suchthat k\in\Re } \), merupakan subruang dari \( \Re^3 \).
Hal ini tetap benar jika $\vec{v}$ adalah vektor nol; dalam kasus tersebut,
garisnya merosot menjadi subruang trivial.
\end{example}

\begin{example}
Rentang himpunan berikut adalah seluruh $\Re^2$.
\begin{equation*}
  \set{\colvec[r]{1 \\ 1},\colvec[r]{1 \\ -1}}
\end{equation*}
Kita mengetahui bahwa rentang tersebut merupakan suatu subruang dari~$\Re^2$.
Untuk memeriksa bahwa rentangnya adalah seluruh~$\Re^2$, kita harus
menunjukkan bahwa setiap anggota $\Re^2$ merupakan kombinasi linear kedua
vektor tersebut.
Jadi, kita bertanya:~bagi vektor mana dengan komponen real $x$ dan~$y$ terdapat
skalar $c_1$ dan $c_2$ yang memenuhi persamaan berikut?
\begin{equation*}
   c_1\colvec[r]{1 \\ 1}+c_2\colvec[r]{1 \\ -1}=\colvec{x \\ y}
  \tag{$*$}
\end{equation*} 
Metode Gauss,
\begin{equation*}
  \begin{linsys}{2}
    c_1  &+  &c_2  &=  &x  \\
    c_1  &-  &c_2  &=  &y
  \end{linsys}
  \grstep{-\rho_1+\rho_2}
  \begin{linsys}{2}
    c_1  &+  &c_2    &=  &x\hfill  \\
         &   &-2c_2  &=  &-x+y
  \end{linsys}
\end{equation*}
bersama dengan substitusi balik, memberikan $c_2=(x-y)/2$ dan $c_1=(x+y)/2$.
Hal ini menunjukkan bahwa bagi setiap $x,y$ terdapat koefisien yang sesuai,
$c_1,c_2$, sehingga~($*$) berlaku\Dash kita dapat menulis setiap anggota
$\Re^2$ sebagai kombinasi linear kedua vektor yang diberikan.
Sebagai contoh, untuk $x=1$ dan $y=2$, koefisien $c_2=-1/2$ dan $c_1=3/2$
memenuhi persamaan tersebut.
\end{example}

Karena rentang merupakan subruang, dan kita mengetahui bahwa cara yang baik
untuk memahami subruang adalah melakukan parametrisasi pada deskripsinya, kita
dapat mencoba memahami rentang suatu himpunan dengan cara tersebut.

\begin{example}
Dalam ruang vektor polinomial kuadrat~\( \polyspace_2 \), tinjau rentang
himpunan \( S=\set{3x-x^2, 2x} \).
Berdasarkan definisi rentang, rentang itu adalah himpunan kombinasi linear
tanpa pembatasan dari keduanya,
$\set{c_1(3x-x^2)+c_2(2x)\suchthat c_1,c_2\in\Re}$.
Jelas bahwa polinomial dalam rentang ini harus memiliki suku konstan nol.
Apakah syarat perlu tersebut juga memadai?

Kita bertanya:~bagi anggota $a_2x^2+a_1x+a_0$ mana dari $\polyspace_2$
terdapat $c_1$ dan $c_2$ sedemikian sehingga
$a_2x^2+a_1x+a_0=c_1(3x-x^2)+c_2(2x)$?
Dua polinomial sama jika koefisiennya sama, sehingga kita menginginkan syarat
pada $a_2$, $a_1$, dan $a_0$ yang menjadikan tripel tersebut solusi sistem
berikut.
\begin{equation*}
  \begin{linsys}{2}
    -c_1  &   &     &=  &a_2   \\
    3c_1  &+  &2c_2 &=  &a_1   \\
          &   &0    &=  &a_0                                   
  \end{linsys}
\end{equation*} 
Metode Gauss dan substitusi balik memberikan $c_1=-a_2$,
$c_2=(3/2)a_2+(1/2)a_1$, dan $0=a_0$.
Jadi,
%the only condition on elements  $a_0+a_1x+a_2x^2$ of the span
%is the condition that we knew: 
selama suku konstannya tidak ada, $a_0=0$, kita dapat memberikan koefisien
$c_1$ dan $c_2$ yang mendeskripsikan polinomial tersebut sebagai anggota
rentang.
Sebagai contoh, untuk polinomial $0-4x+3x^2$, koefisien $c_1=-3$ dan
$c_2=5/2$ memenuhi persamaan.
Jadi, rentang himpunan yang diberikan adalah
$\spanof{S}=\set{a_1x+a_2x^2\suchthat a_1,a_2\in\Re}$. 

Kebetulan, hal ini menunjukkan bahwa himpunan \( \set{x,x^2} \) merentang
subruang yang sama.
Suatu ruang dapat memiliki lebih dari satu himpunan perentang.
Dua himpunan lain yang merentang subruang ini adalah
\( \set{x,x^2,-x+2x^2} \) dan
\( \set{x,x+x^2,x+2x^2,\ldots\,} \).
% (Usually we prefer to work with spanning sets that have only
% a small number of members.)
\end{example}

\begin{example}  \label{ex:SubspRThree}
Gambar di bawah ini menunjukkan subruang-subruang \( \Re^3 \) yang kini kita
ketahui: subruang trivial, garis-garis yang melalui titik asal, bidang-bidang
yang melalui titik asal, dan seluruh ruang.
(Tentu saja, gambar tersebut hanya menampilkan beberapa dari tak hingga banyak
kasus. Ruas garis menghubungkan subhimpunan dengan superhimpunannya.)
Dalam bagian berikutnya, kita akan membuktikan bahwa $\Re^3$ tidak memiliki
jenis subruang lain, sehingga daftar ini sebenarnya lengkap.

Di sini, setiap subruang dideskripsikan sebagai rentang suatu himpunan dengan
jumlah anggota minimal.
Dengan demikian, subruang-subruang tersebut secara alami tersusun dalam
tingkatan\Dash bidang pada satu tingkat, garis pada tingkat lain, dan
seterusnya.
\begin{center}
  \setlength{\unitlength}{4pt}
  \begin{picture}(75,38)(0,-2) %subspaces of R-three
      \thinlines
      \put(45,31){\makebox(0,0)[bl]{
                        \scriptsize \( %\Re^3=
                                   \set{x\colvec[r]{1 \\ 0 \\ 0}
                                              +y\colvec[r]{0 \\ 1 \\ 0}
                                              +z\colvec[r]{0 \\ 0 \\ 1}} \)} }
    %Next the dimension two subspaces
      \put(43,31){\line(-4,-1){25} } % connects R3 to 2,a
      %set 2,a:
      \put(0,22){\makebox(0,0)[l]{\scriptsize\( \set{x\colvec[r]{1 \\ 0 \\ 0}
                                                 +y\colvec[r]{0 \\ 1 \\ 0} }\) }}
      \put(48,29){\line(-3,-1){12} } % connects R3 to 2,b
      %set 2,b:
      \put(20,22){\makebox(0,0)[l]{\scriptsize\( \set{x\colvec[r]{1 \\ 0 \\ 0}
                                                 +z\colvec[r]{0 \\ 0 \\ 1} }\) }}
      \put(53,29){\line(-1,-1){3}} % connects R3 to 2,c
      %set 2,c:
      \put(40,22){\makebox(0,0)[l]{\scriptsize\( \set{x\colvec[r]{1 \\ 1 \\ 0}
                                                 +z\colvec[r]{0 \\ 0 \\ 1} }\) }}
      \put(60,22){\makebox(0,0)[l]{$\cdots$} }
    %Next the dimension one subspaces
      \put(7,18){\line(-1,-4){1} } % connects 2,a to 1,a
      \put(22,18){\line(-3,-1){13} } % connects 2,c to 1,a
      %set 1,a:
      \put(0,10){\makebox(0,0)[l]{\scriptsize\( \set{x\colvec[r]{1 \\ 0 \\ 0}} \)} }
      \put(12,18){\line(1,-2){2} } % connects 2,a to 1,b
      %set 1,b:
      \put(12,10){\makebox(0,0)[l]{\scriptsize\( \set{y\colvec[r]{0 \\ 1 \\ 0}} \)} }
      \put(14,18){\line(2,-1){10} } % connects 2,a to 1,c
      %set 1,c:
      \put(24,10){\makebox(0,0)[l]{\scriptsize\( \set{y\colvec[r]{2 \\ 1 \\ 0}} \)} }
      \put(46,18){\line(-1,-1){3.25} } %connects 2,c to 1,d
      %set 1,d:
      \put(36,10){\makebox(0,0)[l]{\scriptsize\( \set{y\colvec[r]{1 \\ 1 \\ 1}} \)} }
      \put(48,10){\makebox(0,0)[l]{$\cdots$} }
    %Finally, the trivial subspace.
      \put(9,7.5){\line(4,-1){30} } %connects 1,a to trivial
      \put(18.9,7.7){\line(3,-1){20} } %connects 1,b to trivial
      \put(30.4,6.8){\line(2,-1){10} } %connects 1,c to trivial
      \put(41,6){\line(1,-2){1.5} } %connects 1,d to trivial
      \put(45,0){\makebox(0,0){\scriptsize\( \set{\colvec[r]{0 \\ 0 \\ 0} } \)} }
  \end{picture}
\end{center}
\end{example}

Sejauh ini dalam bab ini kita telah melihat bahwa kerangka yang tepat untuk
mempelajari sifat kombinasi linear adalah kumpulan yang tertutup terhadap
kombinasi tersebut.
Dalam subbagian pertama, kita memperkenalkan kumpulan semacam itu, yaitu ruang
vektor, dan melihat beraneka ragam contoh.
Dalam subbagian ini, kita melihat lebih banyak ruang, yakni ruang yang menjadi
subruang dari ruang lain.
Di balik keragaman itu terdapat suatu kesamaan.
\nearbyexample{ex:SubspRThree} di atas menampakkannya:~ruang vektor dan subruang
paling baik dipahami sebagai rentang, khususnya rentang dari sejumlah kecil
vektor. Bagian berikutnya mempelajari himpunan perentang yang minimal.





\begin{exercises}
  \recommended \item
    Manakah dari subhimpunan ruang vektor matriks \( \nbyn{2} \) berikut yang
    merupakan subruang di bawah operasi yang diwarisi?
    Untuk setiap subruang, lakukan parametrisasi pada deskripsinya.
    Untuk setiap himpunan yang bukan subruang, berikan satu syarat yang gagal.
    \begin{exparts}
      \partsitem \( \set{\begin{mat}
                      a  &0  \\
                      0  &b
                    \end{mat}  \suchthat a,b\in\Re}  \)
      \partsitem \( \set{\begin{mat}
                      a  &0  \\
                      0  &b
                    \end{mat}  \suchthat a+b=0} \)
      \partsitem \( \set{\begin{mat}
                      a  &0  \\
                      0  &b
                    \end{mat}  \suchthat a+b=5} \)
      \partsitem \( \set{\begin{mat}
                      a  &c  \\
                      0  &b
                    \end{mat}  \suchthat a+b=0, c\in\Re} \)
    \end{exparts}
    \begin{answer}
      Menurut \nearbylemma{th:SubspIffClosed}, untuk menentukan apakah setiap
      subhimpunan $\matspace_{\nbyn{2}}$ merupakan subruang, kita hanya perlu
      memeriksa apakah subhimpunan tersebut tak kosong dan tertutup.
      \begin{exparts}
        \partsitem Ya, kita dapat dengan mudah memeriksa bahwa himpunan tersebut
          tak kosong dan tertutup. Berikut adalah parametrisasi.
          \begin{equation*}
            \set{a\begin{mat}[r]
                    1  &0  \\
                    0  &0
                  \end{mat}
                 +b\begin{mat}[r]
                    0  &0  \\
                    0  &1  
                   \end{mat}
                 \suchthat a,b\in\Re}
          \end{equation*}
          Parametrisasi tersebut sekaligus menunjukkan bahwa himpunan ini
          merupakan subruang karena dinyatakan sebagai rentang himpunan dua
          matriks, dan setiap rentang merupakan subruang.
        \partsitem Ya; mudah diperiksa bahwa himpunan ini tak kosong dan
          tertutup. Sebagai alternatif, seperti disebutkan dalam jawaban
          sebelumnya, adanya parametrisasi menunjukkan bahwa himpunan ini
          merupakan subruang. Untuk parametrisasi, syarat $a+b=0$ dapat ditulis
          ulang sebagai $a=-b$. Maka kita memperoleh
          \begin{equation*}
            \set{\begin{mat}
                    -b  &0  \\
                    0   &b
                  \end{mat}
                 \suchthat b\in\Re}
            =\set{b\begin{mat}[r]
                    -1  &0  \\
                    0   &1
                  \end{mat}
                 \suchthat b\in\Re}
          \end{equation*}
        \partsitem Tidak. Himpunan ini tidak tertutup terhadap penjumlahan.
          Sebagai contoh,
          \begin{equation*}
            \begin{mat}[r]
              5  &0  \\
              0  &0
            \end{mat}
            +\begin{mat}[r]
              5  &0  \\
              0  &0
            \end{mat}
            =\begin{mat}[r]
              10  &0  \\
              0  &0
            \end{mat}
          \end{equation*}
          bukan anggota himpunan.
          (Himpunan ini juga tidak tertutup terhadap perkalian skalar; misalnya,
          himpunan tersebut tidak memuat matriks nol.)
        \partsitem Ya.
          \begin{equation*}
            \set{b\begin{mat}[r]
                    -1  &0  \\
                    0   &1
                  \end{mat}
                 +c\begin{mat}[r]
                    0  &1  \\
                    0  &0  
                   \end{mat}
                 \suchthat b,c\in\Re}
          \end{equation*}
      \end{exparts}  
     \end{answer}
  \recommended \item 
    Apakah himpunan berikut merupakan subruang dari \( \polyspace_2 \):
    \( \set{a_0+a_1x+a_2x^2\suchthat a_0+2a_1+a_2=4} \)?
    Jika ya, lakukan parametrisasi pada deskripsinya.
    \begin{answer}
      Tidak, himpunan ini tidak tertutup.
      Secara khusus, himpunan ini tidak tertutup terhadap perkalian skalar
      karena tidak memuat polinomial nol.
    \end{answer}
  \item Apakah vektor berikut berada dalam rentang himpunan yang diberikan?
    \begin{equation*}
      \colvec{1 \\ 0 \\ 3}
      \quad
      \set{\colvec{2 \\ 1 \\ -1},
           \colvec{1 \\ -1 \\ 1}}
    \end{equation*}
    \begin{answer}
      Persamaan
      \begin{equation*}
        \colvec{1 \\ 0 \\ 3} = 
                 c_1\colvec{2 \\ 1 \\ -1}
                 +c_2\colvec{1 \\ -1 \\ 1}
      \end{equation*}
      menghasilkan sistem linear
      \begin{equation*}
        \begin{amat}{2}
          2  &1  &1  \\
          1  &-1 &0  \\
          -1 &1  &3
        \end{amat}
        \grstep[(1/2)\rho_1+\rho_3]{(-1/2)\rho_1+\rho_2}
        \begin{amat}{2}
          2  &1    &1  \\
          0  &-3/2 &-1/2  \\
          0  &0  &3
        \end{amat}
      \end{equation*}
      yang tidak memiliki solusi, sehingga vektor tersebut tidak berada dalam
      rentang.
    \end{answer}
  \recommended \item 
    Tentukan apakah vektor tersebut berada dalam rentang himpunan yang
    diberikan di dalam ruangnya.
    \begin{exparts}
      \partsitem \( \colvec[r]{2 \\ 0 \\ 1} \),
        \( \set{\colvec[r]{1 \\ 0 \\ 0},
                \colvec[r]{0 \\ 0 \\ 1}  } \),
        dalam \( \Re^3 \)
      \partsitem \( x-x^3 \),
        \( \set{x^2,2x+x^2,x+x^3} \),
        dalam \( \polyspace_3 \)
      \partsitem \( \begin{mat}[r]
                 0  &1  \\
                 4  &2
               \end{mat}  \),
        \( \set{\begin{mat}[r]
                  1  &0  \\
                  1  &1
                \end{mat},
                \begin{mat}[r]
                  2  &0  \\
                  2  &3
                \end{mat}  } \),
        dalam \( \matspace_{\nbyn{2}} \)
    \end{exparts}
    \begin{answer}
      \begin{exparts}
         \partsitem Ya. Menyelesaikan sistem linear yang berasal dari
           \begin{equation*}
             r_1\colvec[r]{1 \\ 0 \\ 0}+r_2\colvec[r]{0 \\ 0 \\ 1}
               =\colvec[r]{2 \\ 0 \\ 1}
           \end{equation*}
           memberikan \( r_1=2 \) dan \( r_2=1 \).
         \partsitem Ya; sistem linear yang berasal dari
           \( r_1(x^2)+r_2(2x+x^2)+r_3(x+x^3)=x-x^3 \)
           \begin{equation*}
             \begin{linsys}{3}
                   &  &2r_2 &+ &r_3 &= &1  \\
               r_1 &+ &r_2  &  &    &= &0  \\
                   &  &     &  &r_3 &= &-1   
             \end{linsys}
           \end{equation*}
           memberikan \( -1(x^2)+1(2x+x^2)-1(x+x^3)=x-x^3 \).
        \partsitem Tidak; setiap kombinasi kedua matriks yang diberikan
          memiliki entri kanan atas nol.
      \end{exparts}  
    \end{answer}
  \item % \cite{Cleary} credit removed at request Cleary 2022-Sep-13
    Seorang pahlawan super berada di titik asal suatu bidang dua dimensi.
    Pahlawan super itu memiliki dua alat: papan luncur yang dapat bergerak
    sejauh apa pun dalam arah $\binom{3}{1}$ dan karpet ajaib yang dapat
    bergerak sejauh apa pun dalam arah $\binom{1}{2}$.
    \begin{exparts}  
      \partsitem Seorang penjahat bersembunyi di bidang pada titik $(-5,7)$.
        Berapa satuan dengan papan luncur dan karpet ajaib yang harus ditempuh
        pahlawan super untuk mencapai penjahat tersebut?
     \partsitem Adakah tempat di bidang tempat penjahat dapat bersembunyi dengan
        aman tanpa dapat dijangkau?
        Jika ada, berikan satu lokasi semacam itu. Jika tidak, jelaskan
        alasannya.
     \partsitem Pahlawan super dan penjahat dipindahkan ke ruang tiga dimensi,
        dan kini pahlawan super memiliki tiga alat.
        \begin{equation*}
          \text{papan luncur: }\colvec{-1 \\ 0 \\ 3}
          \quad
          \text{karpet ajaib: }\colvec{2 \\ 1 \\ 0}
          \quad
          \text{skuter: }\colvec{5 \\ 4 \\ -9}
        \end{equation*}
        Adakah tempat yang aman bagi penjahat untuk bersembunyi?
        Jika ada, berikan satu lokasi semacam itu; jika tidak, jelaskan
        alasannya.
    \end{exparts}
    \begin{answer}
      \begin{exparts}
        \partsitem Hubungan
          \begin{equation*}
            \colvec{-5 \\ 7}=a\colvec{3 \\ 1}+b\colvec{1 \\ 2}
          \end{equation*}
          memberikan sistem persamaan berikut.
          \begin{equation*}
            \begin{linsys}{2}
              3a &+ &b  &= &-5 \\
               a &+ &2b &= &7
            \end{linsys}
            \grstep{-(1/3)\rho_1+\rho_2}
            \begin{linsys}{2}
              3a &+ &b      &= &-5 \\
                 &  &(5/3)b &= &26/3
            \end{linsys}
          \end{equation*}
          Jadi, $b=26/5$.
          Menyubstitusikan $3a+(26/5)=-5$ memberikan $a=-17/5$.
          Dalam istilah bagian ini, vektor $\binom{-5}{7}$ berada dalam
          rentang himpunan dua vektor lainnya.
        \partsitem Tidak ada tempat untuk bersembunyi.
          Reduksi Metode Gauss berikut,
          \begin{equation*}
            \begin{linsys}{2}
              3a &+ &b  &= &x \\
               a &+ &2b &= &y
            \end{linsys}
            \grstep{-(1/3)\rho_1+\rho_2}
            \begin{linsys}{2}
              3a &+ &b      &= &x\hfill\hbox{} \\
                 &  &(5/3)b &= &-(1/3)x+y
            \end{linsys}
          \end{equation*}
          menunjukkan bahwa bagi setiap $x,y\in\R$ terdapat pasangan
          $a,b\in\R$.
          Dalam istilah bagian ini, rentang kedua vektor adalah seluruh bidang.
        \partsitem Tidak ada tempat untuk bersembunyi.
          Tinjau sistem berikut.
          \begin{equation*}
            a\colvec{-1 \\ 0 \\ 3}+b\colvec{2 \\ 1 \\ 0}+c\colvec{5 \\ 4 \\ -9}
                = \colvec{x \\ y \\ z}
          \end{equation*}
          Setelah reduksi,
          \begin{align*}
            \begin{linsys}{3}
              -a &+ &2b &+ &5c &= &x \\
                 &  &b  &+ &4c &= &y \\
              3a &  &   &- &9c &= &z \\
            \end{linsys}
            & \grstep{3\rho_1+\rho_3}
            \begin{linsys}{3}
              -a &+ &2b  &+ &5c &= &x\hfill\hbox{} \\
                 &  &b   &+ &4c &= &y\hfill\hbox{} \\
                 &  &6b  &+ &6c &= &3x+z \\
            \end{linsys}                                  \\
            & \grstep{-6\rho_2+\rho_3}
            \begin{linsys}{3}
              -a &+ &2b  &+ &5c  &= &x\hfill\hbox{} \\
                 &  &b   &+ &4c  &= &y\hfill\hbox{} \\
                 &  &    & &-18c &= &3x-6y+z \\
            \end{linsys}
          \end{align*}
          terlihat bahwa bagi setiap $x,y,z\in\R$ terdapat $a$, $b$, dan~$c$
          yang diperlukan.
          Ketiga alat pahlawan super merentang seluruh ruang tiga dimensi.
      \end{exparts}
    \end{answer}
  \item 
    Manakah dari fungsi-fungsi berikut yang menjadi anggota rentang
    \( \spanof{\set{\cos^2x,\sin^2x} } \)
    dalam ruang vektor fungsi bernilai real dari satu variabel real?
    \begin{exparts*}
      \partsitem \( f(x)=1 \)
      \partsitem \( f(x)=3+x^2 \)
      \partsitem \( f(x)=\sin x \)
      \partsitem \( f(x)=\cos (2x) \)
    \end{exparts*}
    \begin{answer}
      \begin{exparts}
        \partsitem Ya; fungsi tersebut berada dalam rentang karena
          \( 1\cdot\cos^2x+1\cdot\sin^2x=f(x) \).
        \partsitem Tidak, karena
          \( r_1\cos^2x+r_2\sin^2x=3+x^2 \) tidak memiliki solusi skalar yang
          berlaku bagi semua \( x \).
          Sebagai contoh, menetapkan $x=0$ dan $x=\pi$ memberikan dua persamaan
          $r_1\cdot 1+r_2\cdot 0=3$ dan
          $r_1\cdot 1+r_2\cdot 0=3+\pi^2$, yang saling tak konsisten.
        \partsitem Tidak; perhatikan apa yang terjadi ketika $x=\pi/2$ dan
          $x=3\pi/2$.
        \partsitem Ya, \( \cos (2x)=1\cdot\cos^2(x)-1\cdot\sin^2(x) \).
      \end{exparts}  
     \end{answer}
  \recommended \item 
    Manakah dari himpunan-himpunan berikut yang merentang \( \Re^3 \)?
    Dengan kata lain, himpunan mana yang memiliki sifat bahwa setiap vektor
    bertinggi tiga dapat dinyatakan sebagai kombinasi linear yang sesuai dari
    anggota-anggotanya?
    \begin{exparts*}
      \partsitem \( \set{ \colvec[r]{1 \\ 0 \\ 0},
               \colvec[r]{0 \\ 2 \\ 0},
               \colvec[r]{0 \\ 0 \\ 3}  } \)
      \partsitem \( \set{ \colvec[r]{2 \\ 0 \\ 1},
               \colvec[r]{1 \\ 1 \\ 0},
               \colvec[r]{0 \\ 0 \\ 1}  } \)
      \partsitem \( \set{ \colvec[r]{1 \\ 1 \\ 0},
               \colvec[r]{3 \\ 0 \\ 0}  } \)
      \partsitem \( \set{ \colvec[r]{1 \\ 0 \\ 1},
               \colvec[r]{3 \\ 1 \\ 0},
               \colvec[r]{-1 \\ 0 \\ 0},
               \colvec[r]{2 \\ 1 \\ 5}  } \)
      \partsitem \( \set{ \colvec[r]{2 \\ 1 \\ 1},
               \colvec[r]{3 \\ 0 \\ 1},
               \colvec[r]{5 \\ 1 \\ 2},
               \colvec[r]{6 \\ 0 \\ 2}  } \)
    \end{exparts*}
    \begin{answer}
      \begin{exparts}
         \partsitem Ya, bagi setiap \( x,y,z\in\Re \), persamaan
           \begin{equation*}
              r_1\colvec[r]{1 \\ 0 \\ 0}
              +r_2\colvec[r]{0 \\ 2 \\ 0}
              +r_3\colvec[r]{0 \\ 0 \\ 3}
              =\colvec{x \\ y \\ z}
           \end{equation*}
           memiliki solusi \( r_1=x \), \( r_2=y/2 \), dan
           \( r_3=z/3 \).
         \partsitem Ya, persamaan
           \begin{equation*}
             r_1\colvec[r]{2 \\ 0 \\ 1}
             +r_2\colvec[r]{1 \\ 1 \\ 0}
             +r_3\colvec[r]{0 \\ 0 \\ 1}
             =\colvec{x \\ y \\ z}
           \end{equation*}
           menghasilkan sistem berikut.
           \begin{equation*}
             \begin{linsys}{3}
               2r_1 &+  &r_2  &  &    &=  &x \\
                    &   &r_2  &  &    &=  &y \\
                r_1 &   &     &+ &r_3 &=  &z \\
             \end{linsys}
             \grstep{-(1/2)\rho_1+\rho_3}\repeatedgrstep{(1/2)\rho_2+\rho_3}
             \begin{linsys}{3}
               2r_1 &+  &r_2  &  &    &=  &x\hfill\hbox{} \\
                    &   &r_2  &  &    &=  &y\hfill\hbox{} \\
                    &   &     &  &r_3 &=  &-(1/2)x+(1/2)y+z \\
             \end{linsys}
           \end{equation*}
           Dengan demikian, bagi setiap $x$, $y$, dan $z$, kita dapat
           menghitung \( r_3=(-1/2)x+(1/2)y+z \), \( r_2=y \), dan
           \( r_1=(1/2)x-(1/2)y \).
        \partsitem Tidak. Secara khusus, kita tidak dapat memperoleh vektor
           \begin{equation*}
             \colvec[r]{0 \\ 0 \\ 1}
           \end{equation*}
           sebagai kombinasi linear karena kedua vektor yang diberikan
           memiliki komponen ketiga nol.
       \partsitem Ya. Persamaan
         \begin{equation*}
           r_1\colvec[r]{1 \\ 0 \\ 1}
           +r_2\colvec[r]{3 \\ 1 \\ 0}
           +r_3\colvec[r]{-1\\ 0 \\ 0}
           +r_4\colvec[r]{2 \\ 1 \\ 5}
           =\colvec{x \\ y \\ z}
         \end{equation*}
         menghasilkan reduksi berikut.
         \begin{equation*}
           \begin{amat}{4}
             1  &3  &-1  &2  &x  \\
             0  &1  &0   &1  &y  \\
             1  &0  &0   &5  &z  
           \end{amat}                                              
           \grstep{-\rho_1+\rho_3}\repeatedgrstep{3\rho_2+\rho_3}
           \begin{amat}{4}
             1  &3  &-1  &2  &x \\
             0  &1  &0   &1  &y  \\
             0  &0  &1   &6  &-x+3y+z
           \end{amat}
         \end{equation*}
         Kita memiliki tak hingga banyak solusi.
         Sebagai contoh, kita dapat menetapkan $r_4=0$ lalu menyelesaikan
         $r_3$, $r_2$, dan $r_1$ dalam bentuk $x$, $y$, dan $z$ menggunakan
         metode substitusi balik yang biasa.
       \partsitem Tidak. Persamaan
         \begin{equation*}
           r_1\colvec[r]{2 \\ 1 \\ 1}
           +r_2\colvec[r]{3 \\ 0 \\ 1}
           +r_3\colvec[r]{5 \\ 1 \\ 2}
           +r_4\colvec[r]{6 \\ 0 \\ 2}
           =\colvec{x \\ y \\ z}
         \end{equation*}
         menghasilkan reduksi berikut.
         \begin{multline*}
           \begin{amat}{4}
             2  &3  &5   &6  &x  \\
             1  &0  &1   &0  &y  \\
             1  &1  &2   &2  &z  
           \end{amat}                                             \\
           \grstep[-(1/2)\rho_1+\rho_3]{-(1/2)\rho_1+\rho_2}
           \repeatedgrstep{-(1/3)\rho_2+\rho_3}
           \begin{amat}{4}
             2  &3     &5     &6  &x \\
             0  &-3/2  &-3/2  &-3 &-(1/2)x+y  \\
             0  &0     &0     &0  &-(1/3)x-(1/3)y+z
           \end{amat}
         \end{multline*}
         Hal ini menunjukkan bahwa tidak setiap vektor bertinggi tiga dapat
         dinyatakan demikian.
         Hanya vektor yang memenuhi syarat $-(1/3)x-(1/3)y+z=0$ yang berada
         dalam rentang.
         (Untuk melihat bahwa setiap vektor semacam itu memang dapat
         dinyatakan, tetapkan $r_3=r_4=0$ lalu selesaikan $r_1$ dan $r_2$ dalam
         bentuk $x$, $y$, dan $z$ dengan substitusi balik.)
      \end{exparts}  
    \end{answer}
  \recommended \item 
    Parametrisasikan deskripsi setiap subruang.
    Kemudian nyatakan setiap subruang sebagai suatu rentang.
    \begin{exparts}
      \partsitem  Subhimpunan \( \set{\rowvec{a &b &c}\suchthat a-c=0}   \)
        dari vektor baris berlebar tiga
      \partsitem Subhimpunan dari \( \matspace_{\nbyn{2}} \) berikut
        \begin{equation*}
          \set{\begin{mat}
                 a  &b  \\
                 c  &d
               \end{mat}  \suchthat a+d=0}
        \end{equation*}
      \partsitem Subhimpunan dari \( \matspace_{\nbyn{2}} \) berikut
        \begin{equation*}
          \set{\begin{mat}
                 a  &b  \\
                 c  &d
               \end{mat}  \suchthat \text{\( 2a-c-d=0 \) 
                                                 dan \( a+3b=0 \)} }
        \end{equation*}
      \partsitem Subhimpunan \( \set{a+bx+cx^3\suchthat a-2b+c=0} \) dari
        \( \polyspace_3 \)
      \partsitem Subhimpunan dari \( \polyspace_2 \) yang terdiri atas
        polinomial kuadratik \( p \) sedemikian sehingga \( p(7)=0 \)
    \end{exparts}
    \begin{answer}
      \begin{exparts}
        \partsitem \( \set{\rowvec{c &b &c}\suchthat b,c\in\Re}
                      =\set{b\rowvec{0 &1 &0}+c\rowvec{1 &0 &1}
                        \suchthat b,c\in\Re} \)
           Pilihan yang langsung terlihat untuk himpunan perentangnya ialah
           $\set{\rowvec{0 &1 &0},\rowvec{1 &0 &1}}$.
        \partsitem \( \set{\begin{mat}
                       -d &b  \\
                       c  &d
                      \end{mat} \suchthat b,c,d\in\Re}
                     =\set{b\begin{mat}[r]
                       0  &1  \\
                       0  &0
                      \end{mat}
                     +c\begin{mat}[r]
                       0  &0  \\
                       1  &0
                      \end{mat}
                     +d\begin{mat}[r]
                       -1  &0  \\
                       0  &1
                      \end{mat}  \suchthat b,c,d\in\Re} \)
            Salah satu himpunan yang merentang ruang ini terdiri atas ketiga
            matriks tersebut.
        \partsitem Sistem
          \begin{equation*}
            \begin{linsys}{4}
              a  &+  &3b  &   &   &  &  &=  &0  \\
             2a  &   &    &   &-c &- &d &=  &0  
            \end{linsys}
          \end{equation*}
          menghasilkan \( b=-(c+d)/6 \) dan \( a=(c+d)/2 \).
          Jadi, salah satu deskripsinya adalah sebagai berikut.
          \begin{equation*}
            \set{c\begin{mat}[r]
                       1/2  &-1/6  \\
                       1    &0
                      \end{mat}
                     +d\begin{mat}[r]
                       1/2  &-1/6  \\
                       0    &1
                      \end{mat}  \suchthat c,d\in\Re}
           \end{equation*}
          Hal ini menunjukkan bahwa suatu himpunan perentang bagi subruang ini
          terdiri atas kedua matriks tersebut.
        \partsitem Persamaan $a=2b-c$ menunjukkan bahwa himpunan
           \( \set{(2b-c)+bx+cx^3 \suchthat b,c\in\Re} \)
           sama dengan himpunan
           \( \set{b(2+x)+c(-1+x^3) \suchthat b,c\in\Re}  \).
           Jadi, subruang tersebut merupakan rentang dari himpunan
           $\set{2+x, -1+x^3}$.
        \partsitem Himpunan
          \( \set{a+bx+cx^2\suchthat a+7b+49c=0} \)
          dapat diparametrisasikan sebagai
          \begin{equation*}
            \set{b(-7+x)+c(-49+x^2)\suchthat b,c\in\Re} 
          \end{equation*}
          sehingga memiliki himpunan perentang $\set{-7+x,-49+x^2}$.
      \end{exparts}  
    \end{answer}
  \recommended \item 
    Carilah suatu himpunan yang merentang subruang yang diberikan dari ruang
    yang diberikan.
    (\textit{Petunjuk.}   Parametrisasikan masing-masing.)
    \begin{exparts}
      \partsitem  bidang \( xz \) di \( \Re^3 \)
      \partsitem \( \set{\colvec{x \\ y \\ z}\suchthat 3x+2y+z=0} \)
            di \( \Re^3 \)
      \partsitem \( \set{\colvec{x \\ y \\ z \\ w}\suchthat
                       2x+y+w=0 \text{\ dan\ } y+2z=0} \)
            di \( \Re^4 \)
      \partsitem \( \set{a_0+a_1x+a_2x^2+a_3x^3\suchthat
                        a_0+a_1=0 \text{\ dan\ } a_2-a_3=0} \)
            di \( \polyspace_3 \)
      \partsitem Himpunan \( \polyspace_4 \) di dalam ruang \( \polyspace_4 \)
      \partsitem \( \matspace_{\nbyn{2}} \) di dalam \( \matspace_{\nbyn{2}} \)
    \end{exparts}
    \begin{answer}
      \begin{exparts}
        \partsitem Kita dapat membuat parametrisasi berikut.
          \begin{equation*}
             \set{\colvec{x \\ 0 \\ z}\suchthat x,z\in\Re}
             =\set{x\colvec[r]{1 \\ 0 \\ 0}
                  +z\colvec[r]{0 \\ 0 \\ 1}\suchthat x,z\in\Re}
          \end{equation*}
          Dengan demikian, diperoleh himpunan perentang berikut.
          \begin{equation*}
             \set{\colvec[r]{1 \\ 0 \\ 0},\colvec[r]{0 \\ 0 \\ 1}} 
          \end{equation*}
        \item Berikut adalah suatu parametrisasi beserta himpunan perentang
          yang bersesuaian.
          \begin{equation*}
             \set{y\colvec[r]{-2/3 \\ 1 \\ 0}+z\colvec[r]{-1/3 \\ 0 \\ 1}
                 \suchthat y,z\in\Re }
             \qquad
             \set{\colvec[r]{-2/3 \\ 1 \\ 0},\colvec[r]{-1/3 \\ 0 \\ 1} } 
          \end{equation*}
        \partsitem \( \set{\colvec[r]{1 \\ -2 \\ 1 \\ 0},
                      \colvec[r]{-1/2 \\ 0 \\ 0 \\ 1} } \)
        \partsitem Parametrisasikan deskripsinya sebagai
          \( \set{-a_1+a_1x+a_3x^2+a_3x^3\suchthat a_1,a_3\in\Re } \)
          untuk memperoleh \( \set{-1+x,x^2+x^3}. \)
        \partsitem \( \set{1,x,x^2,x^3,x^4} \)
        \partsitem \( \set{ \begin{mat}[r]
                   1  &0  \\
                   0  &0
                 \end{mat},
                 \begin{mat}[r]
                   0  &1  \\
                   0  &0
                 \end{mat},
                 \begin{mat}[r]
                   0  &0  \\
                   1  &0
                 \end{mat},
                 \begin{mat}[r]
                   0  &0  \\
                   0  &1
                 \end{mat} } \)
      \end{exparts}  
    \end{answer}
  \item 
    Apakah \( \Re^2 \) merupakan subruang dari \( \Re^3 \)?
    \begin{answer}
      Secara teknis, tidak.
      Subruang dari \( \Re^3 \) merupakan himpunan vektor dengan tiga
      komponen, sedangkan \( \Re^2 \) merupakan himpunan vektor dengan dua
      komponen.
      Namun, jelas bahwa \( \Re^2 \) ``sama saja dengan'' subruang berikut dari
      \( \Re^3 \).
      \begin{equation*}
        \set{\colvec{x \\ y \\ 0}\suchthat x,y\in\Re}
      \end{equation*}  
    \end{answer}
  \recommended \item 
    Tentukan apakah masing-masing himpunan berikut merupakan subruang dari
    ruang vektor fungsi bernilai real dengan satu variabel real.
    \begin{exparts}
      \partsitem Fungsi-fungsi
        \definend{genap}\index{fungsi!genap}\index{fungsi genap}
        \( \set{\map{f}{\Re}{\Re} \suchthat f(-x)=f(x) \text{ untuk setiap } x} \).
        Sebagai contoh, dua anggota himpunan ini adalah $f_1(x)=x^2$
        dan $f_2(x)=\cos (x)$.
      \partsitem Fungsi-fungsi
        \definend{ganjil}\index{fungsi!ganjil}\index{fungsi ganjil}
        \( \set{\map{f}{\Re}{\Re} \suchthat f(-x)=-f(x) \text{ untuk setiap } x} \).
        Dua anggotanya adalah $f_3(x)=x^3$ dan $f_4(x)=\sin(x)$.
    \end{exparts}
    \begin{answer}
      Tentu saja, operasi penjumlahan dan perkalian skalarnya diwarisi dari
      ruang yang memuatnya.
      \begin{exparts}
        \partsitem Himpunan ini merupakan subruang.
          Himpunan ini tidak kosong karena memuat sekurang-kurangnya dua fungsi
          contoh yang diberikan.
          Himpunan ini tertutup karena jika \( f_1,f_2 \) genap dan
          \( c_1,c_2 \) adalah skalar, maka berlaku
          \begin{equation*}
            (c_1f_1+c_2f_2)\,(-x)
            =c_1\,f_1(-x)+c_2\,f_2(-x)
            =c_1\,f_1(x)+c_2\,f_2(x)
            =(c_1f_1+c_2f_2)\,(x)
          \end{equation*}
        \partsitem Himpunan ini juga merupakan subruang; pemeriksaannya serupa
          dengan pemeriksaan sebelumnya.
      \end{exparts}  
    \end{answer}
  \item 
    \nearbyexample{ex:SpanSingVec} menyatakan bahwa untuk setiap vektor
    $\vec{v}$ yang merupakan anggota suatu ruang vektor $V$, himpunan
    $\set{r\cdot\vec{v}\suchthat r\in\Re}$ merupakan subruang dari $V$.
    (Himpunan ini hanyalah rentang\index{rentang!himpunan tunggal}
    dari himpunan tunggal $\set{\vec{v}}$.)
    Apakah setiap subruang semacam itu harus merupakan subruang sejati?
    \begin{answer}
      Tidak; subruang itu dapat tidak sejati.
      Sebagai contoh, jika ruang vektornya adalah \( \Re^1 \) dan
      \( \vec{v}\neq\zero\, \), maka subruang
      $\set{r\cdot\vec{v}\suchthat r\in\Re}$ sama dengan seluruh $\Re^1$.
    \end{answer}
  \item 
    Sebuah contoh setelah definisi ruang vektor menunjukkan bahwa himpunan
    solusi suatu sistem linear homogen merupakan ruang vektor.
    Dalam istilah subbagian ini, himpunan tersebut merupakan subruang dari
    $\Re^n$ apabila sistemnya memiliki $n$ variabel.
    Bagaimana dengan sistem linear tak homogen; apakah solusi-solusinya
    membentuk subruang (di bawah operasi yang diwarisi)?
    \begin{answer}
      Tidak, himpunan semacam itu tidak tertutup.
      Salah satu alasannya, himpunan itu tidak memuat vektor nol.
    \end{answer}
  \item \cite{Cleary} 
   Berikan contoh untuk setiap hal berikut atau jelaskan mengapa hal itu
   mustahil dilakukan.
   \begin{exparts}
     \item Subhimpunan tak kosong dari $\matspace_{\nbyn{2}}$ yang bukan
       merupakan subruang.
     \item Himpunan dua vektor di $\Re^2$ yang tidak merentang ruang tersebut.
   \end{exparts}
   \begin{answer}
     \begin{exparts}
       \item Subhimpunan tak kosong dari $\matspace_{\nbyn{2}}$ berikut bukan
         merupakan subruang.
         \begin{equation*}
           A=\set{
             \begin{mat}[r]
               1 &2 \\
               3 &4
             \end{mat},
             \begin{mat}[r]
               5 &6 \\
               7 &8
             \end{mat}}
         \end{equation*}
         Salah satu alasan bahwa himpunan ini bukan subruang dari
         $\matspace_{\nbyn{2}}$ adalah karena himpunan ini tidak memuat matriks
         nol.
         (Alasan lainnya adalah bahwa himpunan ini tidak tertutup terhadap
         penjumlahan, karena jumlah kedua matriks itu bukan anggota $A$.
         Himpunan ini juga tidak tertutup terhadap perkalian skalar.)
        \item Himpunan dua vektor berikut tidak merentang $\Re^2$.
          \begin{equation*}
            \set{\colvec[r]{1 \\ 1},\colvec[r]{3 \\ 3}}
          \end{equation*}
          Tidak ada kombinasi linear dari kedua vektor ini yang dapat
          menghasilkan vektor dengan komponen kedua yang tidak sama dengan
          komponen pertamanya.
      \end{exparts}
   \end{answer}
  \item 
    \nearbyexample{ex:SubspRThree} menunjukkan bahwa
    $\Re^3$ memiliki tak hingga banyak subruang.
    Apakah setiap ruang nontrivial memiliki tak hingga banyak subruang?
    \begin{answer}
      Tidak.
      Satu-satunya subruang dari \( \Re^1 \) adalah ruang itu sendiri dan
      subruang trivialnya.
      Setiap subruang $S$ dari $\Re$ yang memuat anggota tak nol $\vec{v}$
      harus memuat himpunan semua kelipatan skalarnya,
      $\set{r\cdot\vec{v}\suchthat r\in\Re}$.
      Akan tetapi, himpunan ini sama dengan seluruh $\Re$.
    \end{answer}
  \item \label{exer:SubspIffClosed}
    Selesaikan bukti \nearbylemma{th:SubspIffClosed}.
    \begin{answer}
      Butir~(1) diperiksa dalam teks.

      Butir~(2) memiliki lima syarat.
      Pertama, untuk ketertutupan, jika \( c\in\Re \) dan \( \vec{s}\in S \),
      maka \( c\cdot\vec{s}\in S \) karena
      \( c\cdot\vec{s}=c\cdot\vec{s}+0\cdot\zero \).
      Kedua, karena operasi pada \( S \) diwarisi dari \( V \), untuk
      \( c,d\in\Re \) dan \( \vec{s}\in S \), hasil kali skalar
      \( (c+d)\cdot\vec{s}\, \) di \( S \) sama dengan hasil kali
      \( (c+d)\cdot\vec{s}\, \) di \( V \), dan hasil itu sama dengan
      \( c\cdot\vec{s}+d\cdot\vec{s}\, \) di \( V \), yang sama dengan
      \( c\cdot\vec{s}+d\cdot\vec{s}\, \) di \( S \).

      Pemeriksaan syarat ketiga, keempat, dan kelima serupa dengan pemeriksaan
      syarat kedua yang baru saja diberikan.
    \end{answer}
  \item 
    Tunjukkan bahwa setiap ruang vektor hanya memiliki satu subruang trivial.
    \begin{answer}
      Sebuah latihan pada subbagian sebelumnya menunjukkan bahwa setiap ruang
      vektor hanya memiliki satu vektor nol (yakni, hanya ada satu vektor yang
      menjadi unsur identitas aditif ruang tersebut).
      Akan tetapi, ruang trivial hanya memiliki satu unsur dan unsur itu harus
      merupakan vektor nol (tunggal) ini.
    \end{answer}
  \item 
    Tunjukkan bahwa untuk setiap subhimpunan \( S \) dari suatu ruang vektor,
    rentang dari rentang sama dengan rentang itu sendiri,
    \( \spanof{ \spanof{S} }=\spanof{S} \).
    (\textit{Petunjuk.}
    Anggota $\spanof{S}$ adalah kombinasi linear dari anggota $S$.
    Anggota $\spanof{\spanof{S}}$ adalah kombinasi linear dari kombinasi
    linear anggota $S$.)
    \begin{answer}
      Seperti yang disarankan oleh petunjuk, alasan dasarnya adalah Lema
      Kombinasi Linear dari bab pertama.
      Untuk bukti lengkap, kita akan menunjukkan bahwa kedua himpunan saling
      memuat.

      Pemuatan pertama, \( \spanof{ \spanof{S} }\supseteq\spanof{S}  \),
      merupakan contoh dari fakta yang lebih umum dan jelas bahwa untuk setiap
      subhimpunan \( T \) dari suatu ruang vektor, \( \spanof{T}\supseteq T \).

      Untuk pemuatan sebaliknya, yaitu
      \( \spanof{ \spanof{S} }\subseteq\spanof{S} \), ambil $m$ vektor dari
      \( \spanof{S} \), yakni
      \( c_{1,1}\vec{s}_{1,1}+\cdots+c_{1,n_1}\vec{s}_{1,n_1} \), \ldots,
      \( c_{m,1}\vec{s}_{m,1}+\cdots+c_{m,n_m}\vec{s}_{m,n_m} \),
      dan perhatikan bahwa setiap kombinasi linear dari vektor-vektor tersebut,
      \begin{equation*}
         r_1(c_{1,1}\vec{s}_{1,1}+\cdots+c_{1,n_1}\vec{s}_{1,n_1})+\cdots
        +r_m(c_{m,1}\vec{s}_{m,1}+\cdots+c_{m,n_m}\vec{s}_{m,n_m})
      \end{equation*}
      merupakan kombinasi linear dari unsur-unsur \( S \),
      \begin{equation*}
        = (r_1c_{1,1})\vec{s}_{1,1}+\cdots+(r_1c_{1,n_1})\vec{s}_{1,n_1}+\cdots
        +(r_mc_{m,1})\vec{s}_{m,1}+\cdots+(r_mc_{m,n_m})\vec{s}_{m,n_m}
      \end{equation*}
      sehingga berada dalam \( \spanof{S} \).
      Dengan kata lain, cukup ingat bahwa suatu kombinasi linear dari
      kombinasi-kombinasi linear (anggota $S$) merupakan kombinasi linear
      (sekali lagi, dari anggota $S$).
    \end{answer}
  \item 
    Semua subruang yang telah kita lihat menggunakan nol dalam deskripsinya
    dengan satu atau lain cara.
    Sebagai contoh, subruang dalam \nearbyexample{ex:SubspacesRTwo} terdiri
    atas semua vektor dari $\Re^2$ yang komponen keduanya nol.
    Sebaliknya, kumpulan vektor dari $\Re^2$ yang komponen keduanya satu tidak
    membentuk subruang (kumpulan itu tidak tertutup terhadap perkalian skalar).
    Contoh lainnya adalah \nearbyexample{ex:PlaneSubspRThree}, yang syarat pada
    vektornya adalah bahwa jumlah ketiga komponennya nol.
    Jika syarat di sana menyatakan bahwa jumlah ketiga komponennya satu, maka
    himpunan itu bukan subruang (sekali lagi, ketertutupan tidak terpenuhi).
    Namun, ketergantungan pada nol tidak mutlak diperlukan.
    Pertimbangkan himpunan
    \begin{equation*}
       \set{\colvec{x \\ y \\ z}\suchthat x+y+z=1}
    \end{equation*}
    dengan operasi-operasi berikut.
    \begin{equation*}
       \colvec{x_1 \\ y_1 \\ z_1}+\colvec{x_2 \\ y_2 \\ z_2}
       =\colvec{x_1+x_2-1 \\ y_1+y_2 \\ z_1+z_2}
       \qquad
       r\colvec{x \\ y \\ z}=\colvec{rx-r+1 \\ ry \\ rz}
    \end{equation*}
    \begin{exparts}
      \partsitem Tunjukkan bahwa himpunan itu bukan subruang dari $\Re^3$.
         (\textit{Petunjuk.} Lihat \nearbyexample{ex:OperNotInherit}).
      \partsitem Tunjukkan bahwa himpunan itu merupakan ruang vektor.
         Perhatikan bahwa, berdasarkan butir sebelumnya,
         \nearbylemma{th:SubspIffClosed} tidak dapat diterapkan.
      \partsitem Tunjukkan bahwa setiap subruang dari $\Re^3$ harus melalui
         titik asal, sehingga setiap subruang dari $\Re^3$ harus melibatkan nol
         dalam deskripsinya.
         Apakah kebalikannya berlaku?
          Apakah setiap subhimpunan dari $\Re^3$ yang memuat titik asal menjadi
          subruang jika diberi operasi yang diwarisi?
    \end{exparts}
    \begin{answer}
      \begin{exparts}
         \partsitem Himpunan itu bukan subruang karena operasi-operasi ini bukan
           operasi yang diwarisi.
           Salah satu alasannya, di ruang ini,
           \begin{equation*}
             0\cdot\colvec{x \\ y \\ z}=\colvec[r]{1 \\ 0 \\ 0}
           \end{equation*}
           sedangkan hal ini tentu saja tidak berlaku di $\Re^3$.
         \partsitem Kita dapat menggabungkan argumen yang menunjukkan
           ketertutupan terhadap penjumlahan dengan argumen yang menunjukkan
           ketertutupan terhadap perkalian skalar menjadi satu argumen yang
           menunjukkan ketertutupan terhadap kombinasi linear dua vektor.
           Jika $r_1,r_2,x_1,x_2,y_1,y_2,z_1,z_2$ berada di $\Re$, maka
           \begin{multline*}
              r_1\colvec{x_1 \\ y_1 \\ z_1}
              +r_2\colvec{x_2 \\ y_2 \\ z_2}
              =\colvec{r_1x_1-r_1+1 \\ r_1y_1 \\ r_1z_1}
               +\colvec{r_2x_2-r_2+1 \\ r_2y_2 \\ r_2z_2}                 \\
            =\colvec{r_1x_1-r_1+r_2x_2-r_2+1 \\ r_1y_1+r_2y_2 \\ r_1z_1+r_2z_2}
           \end{multline*} 
           (perhatikan bahwa definisi penjumlahan di ruang ini menggabungkan
           komponen pertama sebagai $(r_1x_1-r_1+1)+(r_2x_2-r_2+1)-1$,
           sehingga komponen pertama vektor terakhir tidak berakhir dengan
           `$\hbox{}+2$').
           Menjumlahkan ketiga komponen vektor terakhir menghasilkan
           $r_1(x_1-1+y_1+z_1)+r_2(x_2-1+y_2+z_2)+1=r_1\cdot0+r_2\cdot0+1=1$.

           Sebagian besar pemeriksaan syarat lainnya mudah (meskipun
           ketidaklaziman operasinya membuat pemeriksaan itu tidak rutin).
           Komutativitas penjumlahan diperiksa sebagai berikut.
           \begin{equation*}
             \colvec{x_1 \\ y_1 \\ z_1}+\colvec{x_2 \\ y_2 \\ z_2}
             =\colvec{x_1+x_2-1 \\ y_1+y_2 \\ z_1+z_2}
             =\colvec{x_2+x_1-1 \\ y_2+y_1 \\ z_2+z_1}
             =\colvec{x_2 \\ y_2 \\ z_2}+\colvec{x_1 \\ y_1 \\ z_1}
           \end{equation*}
           Untuk asosiativitas penjumlahan, kita memiliki
           \begin{equation*}
             (\colvec{x_1 \\ y_1 \\ z_1}+\colvec{x_2 \\ y_2 \\ z_2})
              +\colvec{x_3 \\ y_3 \\ z_3}
             =\colvec{(x_1+x_2-1)+x_3-1 \\ (y_1+y_2)+y_3 \\ (z_1+z_2)+z_3}
           \end{equation*}
           sedangkan
           \begin{equation*}
             \colvec{x_1 \\ y_1 \\ z_1}
             +(\colvec{x_2 \\ y_2 \\ z_2}+\colvec{x_3 \\ y_3 \\ z_3})
             =\colvec{x_1+(x_2+x_3-1)-1 \\ y_1+(y_2+y_3) \\ z_1+(z_2+z_3)}
           \end{equation*}
           dan keduanya sama.
           Unsur identitas terhadap operasi penjumlahan ini bekerja sebagai
           berikut:
           \begin{equation*}
             \colvec{x \\ y \\ z}+\colvec[r]{1 \\ 0 \\ 0}
             =\colvec{x+1-1 \\ y+0 \\ z+0}
             =\colvec{x \\ y \\ z}
           \end{equation*}
           dan invers aditifnya diperiksa dengan cara serupa.
           \begin{equation*}
             \colvec{x \\ y \\ z}+\colvec{-x+2 \\ -y \\ -z}
             =\colvec{x+(-x+2)-1 \\ y-y \\ z-z}
             =\colvec[r]{1 \\ 0 \\ 0}
           \end{equation*}

           Syarat-syarat perkalian skalar juga mudah diperiksa.
           Untuk syarat pertama,
           \begin{equation*}
             (r+s)\colvec{x \\ y \\ z}
             =\colvec{(r+s)x-(r+s)+1 \\ (r+s)y \\ (r+s)z}
           \end{equation*}
           sedangkan
           \begin{multline*}
             r\colvec{x \\ y \\ z}+s\colvec{x \\ y \\ z}
             =\colvec{rx-r+1 \\ ry \\ rz}+\colvec{sx-s+1 \\ sy \\ sz}  \\
             =\colvec{(rx-r+1)+(sx-s+1)-1 \\ ry+sy \\ rz+sz}
           \end{multline*}
           dan keduanya sama.
           Syarat kedua membandingkan
           \begin{equation*}
             r\cdot(\colvec{x_1 \\ y_1 \\ z_1}+\colvec{x_2 \\ y_2 \\ z_2})
             =r\cdot\colvec{x_1+x_2-1 \\ y_1+y_2 \\ z_1+z_2}
             =\colvec{r(x_1+x_2-1)-r+1 \\ r(y_1+y_2) \\ r(z_1+z_2)}
           \end{equation*}
           dengan
           \begin{multline*}
             r\colvec{x_1 \\ y_1 \\ z_1}+r\colvec{x_2 \\ y_2 \\ z_2}
             =\colvec{rx_1-r+1 \\ ry_1 \\ rz_1}
                     +\colvec{rx_2-r+1 \\ ry_2 \\ rz_2}                  \\
             =\colvec{(rx_1-r+1)+(rx_2-r+1)-1 \\ ry_1+ry_2 \\ rz_1+rz_2}
           \end{multline*}
           dan keduanya sama.
           Untuk syarat ketiga,
           \begin{equation*}
             (rs)\colvec{x \\ y \\ z}
             =\colvec{rsx-rs+1 \\ rsy \\ rsz}
           \end{equation*}
           sedangkan
           \begin{equation*}
             r(s\colvec{x \\ y \\ z})
             =r(\colvec{sx-s+1 \\ sy \\ sz})
             =\colvec{r(sx-s+1)-r+1 \\ rsy \\ rsz}
           \end{equation*}
           dan keduanya sama.
           Untuk perkalian skalar dengan $1$, kita memperoleh
           \begin{equation*}
             1\cdot\colvec{x \\ y \\ z}
             =\colvec{1x-1+1 \\ 1y \\ 1z}
             =\colvec{x \\ y \\ z}
           \end{equation*}
           Dengan demikian, kedua operasi ini memenuhi semua syarat ruang
           vektor.

           \textit{Catatan.}
           Salah satu cara memahami ruang vektor ini adalah memandangnya
           sebagai bidang di $\Re^3$
           \begin{equation*}
             P=\set{\colvec{x \\ y \\ z}\suchthat x+y+z=0}
           \end{equation*}
           yang digeser sejauh $1$ dari titik asal sepanjang sumbu-$x$.
           Dengan demikian, penjumlahannya menjadi sebagai berikut: untuk
           menjumlahkan dua anggota ruang ini,
           \begin{equation*}
             \colvec{x_1 \\ y_1 \\ z_1},\;\colvec{x_2 \\ y_2 \\ z_2}
           \end{equation*}
           (dengan $x_1+y_1+z_1=1$ dan $x_2+y_2+z_2=1$), geser keduanya kembali
           sejauh $1$ agar berada di $P$, lalu jumlahkan seperti biasa,
           \begin{equation*}
             \colvec{x_1-1 \\ y_1 \\ z_1}+\colvec{x_2-1 \\ y_2 \\ z_2}
             =\colvec{x_1+x_2-2 \\ y_1+y_2 \\ z_1+z_2}
             \qquad\text{(di $P$)}
           \end{equation*}
           kemudian geser hasilnya kembali sejauh $1$ sepanjang sumbu-$x$.
           \begin{equation*}
             \colvec{x_1+x_2-1 \\ y_1+y_2 \\ z_1+z_2}.
           \end{equation*}
           Perkalian skalarnya serupa.
         \partsitem Agar subruang itu tertutup terhadap perkalian skalar yang
           diwarisi, untuk setiap anggota $\vec{v}$ dari subruang tersebut,
           \begin{equation*}
             0\cdot\vec{v}=\colvec[r]{0 \\ 0 \\ 0}
           \end{equation*}
           juga harus menjadi anggotanya.

           Kebalikannya tidak berlaku.
           Berikut adalah subhimpunan dari $\Re^3$ yang memuat titik asal,
           \begin{equation*}
             \set{\colvec[r]{0 \\ 0 \\ 0},\colvec[r]{1 \\ 0 \\ 0}}
           \end{equation*}
           (subhimpunan ini hanya memiliki dua unsur) tetapi bukan merupakan
           subruang.
      \end{exparts}
    \end{answer}
  \item 
    Kita dapat memberikan pembenaran bagi konvensi bahwa jumlah nol buah
    vektor sama dengan vektor nol.
    Pertimbangkan jumlah tiga vektor
    $\vec{v}_1+\vec{v}_2+\vec{v}_3$ berikut.
    \begin{exparts}
      \partsitem Berapakah selisih antara jumlah tiga vektor ini dan jumlah dua
        vektor pertama di antaranya?
      \partsitem Berapakah selisih antara jumlah sebelumnya dan jumlah yang
        hanya terdiri atas vektor pertama?
      \partsitem Berapakah seharusnya selisih antara jumlah satu vektor
        sebelumnya dan jumlah nol buah vektor?
      \partsitem Jadi, bagaimana seharusnya jumlah nol buah vektor didefinisikan?
    \end{exparts}
    \begin{answer}
      \begin{exparts}
        \item $(\vec{v}_1+\vec{v}_2+\vec{v}_3)-(\vec{v}_1+\vec{v}_2)
                 =\vec{v}_3$
        \item $(\vec{v}_1+\vec{v}_2)-(\vec{v}_1)
                 =\vec{v}_2$
        \item Tentunya, $\vec{v}_1$.
        \item Mengambil jumlah satu vektor itu lalu mengurangkan menghasilkan
          ($\vec{v}_1)-\vec{v}_1=\zero$.
      \end{exparts}
    \end{answer}
  \item 
    Apakah suatu ruang ditentukan oleh subruang-subruangnya?
    Artinya, jika dua ruang vektor memiliki subruang-subruang yang sama, apakah
    kedua ruang itu harus sama?
    \begin{answer}
      Ya; setiap ruang merupakan subruang dari dirinya sendiri, sehingga
      masing-masing ruang memuat ruang yang lain.
    \end{answer}
  \item 
     \begin{exparts}
      \partsitem Berikan himpunan yang tertutup terhadap perkalian skalar tetapi
        tidak terhadap penjumlahan.
      \partsitem Berikan himpunan yang tertutup terhadap penjumlahan tetapi tidak
        terhadap perkalian skalar.
      \partsitem Berikan himpunan yang tidak tertutup terhadap keduanya.
    \end{exparts}
    \begin{answer}
      \begin{exparts}
        \partsitem Gabungan sumbu-\( x \) dan sumbu-\( y \) di \( \Re^2 \) merupakan
          salah satu contohnya.
        \partsitem Himpunan bilangan bulat, sebagai subhimpunan dari
          \( \Re^1 \), merupakan salah satu contohnya.
        \partsitem Subhimpunan \( \set{\vec{v}} \) dari \( \Re^2 \) merupakan
          salah satu contohnya, dengan $\vec{v}$ sebarang vektor tak nol.
      \end{exparts}  
     \end{answer}
  \item 
    Tunjukkan bahwa rentang suatu himpunan vektor tidak bergantung pada urutan
    vektor-vektor tersebut dicantumkan dalam himpunan itu.
    \begin{answer}
      Karena penjumlahan ruang vektor bersifat komutatif, pengurutan ulang
      suku-suku jumlah tidak mengubah suatu kombinasi linear.
    \end{answer}
  \item  
    Subruang trivial manakah yang merupakan rentang dari himpunan kosong?
    Apakah subruang itu
    \begin{equation*}
      \set{\colvec[r]{0 \\ 0 \\ 0}}\subseteq \Re^3,
      \quad\text{atau}\quad
      \set{0+0x}\subseteq \polyspace_1,
    \end{equation*}
    atau suatu subruang lain?
    \begin{answer}
      Kita selalu mempertimbangkan rentang tersebut dalam konteks suatu ruang
      yang memuatnya.
    \end{answer}
  \item   
    Tunjukkan bahwa jika suatu vektor berada dalam rentang sebuah himpunan,
    maka menambahkan vektor tersebut ke himpunan itu tidak akan memperbesar
    rentangnya.
    Apakah pernyataan tersebut juga berlaku ``hanya jika''?
    \begin{answer}
      Pernyataan itu berlaku dalam kedua arah, ``jika'' dan ``hanya jika''.

      Untuk arah ``jika'', misalkan \( S \) adalah subhimpunan suatu ruang
      vektor \( V \) dan andaikan \( \vec{v}\in\spanof{S} \), sehingga
      \( \vec{v}=c_1\vec{s}_1+\dots+c_n\vec{s}_n \), dengan
      \( c_1,\ldots,c_n \) adalah skalar dan
      \( \vec{s}_1,\ldots,\vec{s}_n\in S \).
      Kita harus menunjukkan bahwa
      \( \spanof{S\union\set{\vec{v}} }=\spanof{S} \).

      Pemuatan dalam satu arah,
      \( \spanof{S}\subseteq\spanof{S\union\set{\vec{v}} } \), sudah jelas.
      Untuk arah sebaliknya,
      \( \spanof{S\union\set{\vec{v}} }\subseteq\spanof{S} \), perhatikan
      bahwa jika suatu vektor berada dalam himpunan di ruas kiri, maka vektor
      itu berbentuk
      \( d_0\vec{v}+d_1\vec{t}_1+\dots+d_m\vec{t}_m \), dengan semua \( d \)
      merupakan skalar dan semua \( \vec{t}\, \) berada dalam \( S \).
      Tulis ulang bentuk tersebut sebagai
      \( d_0(c_1\vec{s}_1+\dots+c_n\vec{s}_n)
      +d_1\vec{t}_1+\cdots+d_m\vec{t}_m \), lalu perhatikan bahwa hasilnya
      merupakan anggota rentang \( S \).

      Untuk arah ``hanya jika'', \( \vec{v}\in
      \spanof{S\union\set{\vec{v}}} \). Jadi, jika penambahan \( \vec{v} \)
      tidak mengubah rentang, maka \( \vec{v}\in\spanof{S} \).
    \end{answer}
  \recommended \item
    Subruang merupakan subhimpunan, sehingga wajar jika kita mempertimbangkan
    bagaimana hubungan ``merupakan subruang dari'' berinteraksi dengan operasi
    himpunan yang biasa.
    \begin{exparts}
      \partsitem Jika \( A,B \) merupakan subruang dari suatu ruang vektor,
        apakah irisan keduanya, \( A\intersection B \), harus merupakan
        subruang?
        Selalu? Kadang-kadang? Tidak pernah?
      \partsitem Apakah gabungan \( A\union B \) harus merupakan subruang?
      \partsitem Jika \( A \) merupakan subruang dari suatu~$V$, apakah
        komplemen himpunannya, $V-A$, harus merupakan subruang?
    \end{exparts}
    (\textit{Petunjuk.} Cobalah beberapa subruang uji dari
    \nearbyexample{ex:SubspRThree}.)
    \begin{answer}
      \begin{exparts}
        \partsitem Selalu.

          Andaikan \( A,B \) merupakan subruang dari \( V \).
          Perhatikan bahwa irisannya tidak kosong karena keduanya memuat vektor
          nol.
          Jika \( \vec{v},\vec{w}\in A\intersection B \) dan \( r,s \) adalah
          skalar, maka \( r\vec{v}+s\vec{w}\in A \) karena kedua vektor tersebut
          berada dalam \( A \), sehingga kombinasi linearnya berada dalam \( A \), dan
          \(r\vec{v}+s\vec{w}\in B \) karena alasan yang sama.
          Dengan demikian, irisan tersebut tertutup.
          Sekarang \nearbylemma{th:SubspIffClosed} dapat diterapkan.
        \partsitem Kadang-kadang (lebih tepatnya, hanya jika
          \( A\subseteq B \) atau \( B\subseteq A \)).

          Untuk melihat bahwa jawabannya bukan ``selalu'', ambil
          \( V \) sebagai \( \Re^3 \), ambil \( A \) sebagai sumbu-$x$, dan
          \( B \) sebagai sumbu-\( y \).
          Perhatikan bahwa
          \begin{equation*}
            \colvec[r]{1 \\ 0 \\ 0}\in A \text{ dan }\colvec[r]{0 \\ 1 \\ 0}\in B
            \quad\text{tetapi}\quad
            \colvec[r]{1 \\ 0 \\ 0}+\colvec[r]{0 \\ 1 \\ 0}\not\in A\union B
          \end{equation*}
          karena jumlahnya tidak berada dalam \( A \) maupun \( B \).

          Jawabannya bukan ``tidak pernah'' karena jika \( A\subseteq B \) atau
          \( B\subseteq A \), maka jelas bahwa \( A\union B \) merupakan
          subruang.

          Untuk menunjukkan bahwa \( A\union B \) merupakan subruang hanya jika
          salah satu subruang memuat subruang lainnya, andaikan
          \( A\not\subseteq B \) dan \( B\not\subseteq A \), lalu buktikan bahwa
          gabungannya bukan subruang.
          Asumsi bahwa \( A \) bukan subhimpunan dari \( B \) berarti ada
          \( \vec{a}\in A \) dengan \( \vec{a}\not\in B \).
          Asumsi lainnya memberikan \( \vec{b}\in B \) dengan
          \( \vec{b}\not\in A \).
          Pertimbangkan \( \vec{a}+\vec{b} \).
          Perhatikan bahwa jumlah ini bukan anggota \( A \); sebab jika
          sebaliknya, maka \( (\vec{a}+\vec{b})-\vec{a} \) akan berada dalam
          \( A \), padahal tidak demikian.
          Demikian pula, jumlah tersebut bukan anggota \( B \).
          Jadi, jumlah itu bukan anggota \( A\union B \), sehingga gabungannya
          bukan subruang.
        \partsitem Tidak pernah.
          Karena \( A \) merupakan subruang, subruang itu memuat vektor nol;
          akibatnya, komplemen $A$, yaitu $V-A$, tidak memuat vektor nol.
          Tanpa vektor nol, komplemen tersebut tidak dapat menjadi ruang vektor.
      \end{exparts}  
    \end{answer}
  \item 
    Apakah rentang suatu himpunan bergantung pada ruang yang memuatnya?
    Artinya, jika \( W \) merupakan subruang dari \( V \) dan \( S \) adalah
    subhimpunan dari \( W \) (sehingga juga merupakan subhimpunan dari
    \( V \)), mungkinkah rentang \( S \) dalam \( W \) berbeda dari rentang
    \( S \) dalam \( V \)?
    \begin{answer}
      Rentang suatu himpunan tidak bergantung pada ruang yang memuatnya.
      Kombinasi linear vektor-vektor dari \( S \) menghasilkan jumlah yang sama,
      baik kita menganggap operasinya sebagai operasi \( W \) maupun sebagai
      operasi \( V \), karena operasi \( W \) diwarisi dari \( V \).
    \end{answer}
  \item 
    Apakah relasi ``merupakan subruang dari'' bersifat transitif?
    Artinya, jika $V$ merupakan subruang dari $W$ dan $W$ merupakan subruang
    dari $X$, apakah $V$ harus merupakan subruang dari $X$?
    \begin{answer}
      Ya; terapkan \nearbylemma{th:SubspIffClosed}.
      (Anda perlu mempertimbangkan hal berikut.
      Misalkan \( B \) merupakan subruang dari suatu ruang vektor \( V \) dan
      misalkan \( A\subseteq B\subseteq V \) merupakan subruang.
      Dari ruang manakah \( A \) mewarisi operasinya?
      Jawabannya tidak menjadi masalah\Dash \( A \) akan mewarisi operasi yang
      sama dalam kedua kasus.)
    \end{answer}
  \item
    Karena ``rentang dari'' merupakan operasi pada himpunan, wajar jika kita
    mempertimbangkan bagaimana operasi itu berinteraksi dengan operasi himpunan
    yang biasa.
    \begin{exparts}
      \partsitem Jika \( S\subseteq T \) merupakan subhimpunan suatu ruang
        vektor, apakah
        \( \spanof{S}\subseteq\spanof{T} \)?
        Selalu? Kadang-kadang? Tidak pernah?
      \partsitem Jika \( S,T \) merupakan subhimpunan suatu ruang vektor, apakah
        \( \spanof{S\union T}=\spanof{S}\union\spanof{T} \)?
      \partsitem Jika \( S,T \) merupakan subhimpunan suatu ruang vektor, apakah
        \( \spanof{S\intersection T}=\spanof{S}\intersection\spanof{T} \)?
      \partsitem Apakah rentang komplemen sama dengan komplemen rentang?
    \end{exparts}
    \begin{answer}
      \begin{exparts}
         \partsitem Selalu;
           jika \( S\subseteq T \), maka kombinasi linear unsur-unsur \( S \)
           juga merupakan kombinasi linear unsur-unsur \( T \).
         \partsitem Kadang-kadang (lebih tepatnya, jika dan hanya jika
           \( \spanof{S}\subseteq\spanof{T} \) atau
           \( \spanof{T}\subseteq\spanof{S} \)).

           Jawabannya bukan ``selalu'', seperti ditunjukkan oleh contoh berikut
           dari \( \Re^3 \):
           \begin{equation*}
             S=\set{\colvec[r]{1 \\ 0 \\ 0},\colvec[r]{0 \\ 1 \\ 0}},\quad
             T=\set{\colvec[r]{1 \\ 0 \\ 0},\colvec[r]{0 \\ 0 \\ 1}}
           \end{equation*}
           karena
           \begin{equation*}
             \colvec[r]{1 \\ 1 \\ 1}\in\spanof{S\union T}
             \qquad
             \colvec[r]{1 \\ 1 \\ 1}\not\in\spanof{S}\union \spanof{T}
           \end{equation*}

           Jawabannya bukan ``tidak pernah'' karena jika rentang salah satu
           himpunan memuat rentang himpunan lainnya, kesamaannya jelas.
           Untuk membuktikan bahwa syarat tersebut juga perlu, andaikan
           \( \spanof{S}\not\subseteq\spanof{T} \) dan
           \( \spanof{T}\not\subseteq\spanof{S} \). Dengan demikian, terdapat
           \( \vec{s}\in S \) dengan \( \vec{s}\notin\spanof{T} \), serta
           \( \vec{t}\in T \) dengan \( \vec{t}\notin\spanof{S} \).
           Vektor \( \vec{s}+\vec{t}\in\spanof{S\union T} \), tetapi tidak
           berada dalam \( \spanof{S} \), sebab jika berada di sana maka
           \( \vec{t}=(\vec{s}+\vec{t})-\vec{s}\in\spanof{S} \). Dengan
           alasan yang sama, vektor itu tidak berada dalam \( \spanof{T} \).
           Jadi \( \vec{s}+\vec{t}\notin\spanof{S}\union\spanof{T} \).
         \partsitem Kadang-kadang.

           Jelas bahwa
           \( \spanof{S\intersection T}
              \subseteq\spanof{S}\intersection\spanof{T} \)
           karena setiap kombinasi linear vektor-vektor dari
           \( S\intersection T \) merupakan kombinasi vektor-vektor dari
           \( S \) dan juga kombinasi vektor-vektor dari \( T \).

           Pemuatan dalam arah sebaliknya tidak selalu berlaku.
           Sebagai contoh, di \( \Re^2 \), ambil
           \begin{equation*}
             S=\set{\colvec[r]{1 \\ 0},\colvec[r]{0 \\ 1}},\quad
             T=\set{\colvec[r]{2 \\ 0}}
           \end{equation*}
           sehingga \( \spanof{S}\intersection\spanof{T} \) merupakan
           sumbu-\( x \), sedangkan \( \spanof{S\intersection T} \) merupakan
           subruang trivial.

           Mencirikan dengan tepat kapan kesamaan berlaku tidaklah mudah.
           Jelas bahwa kesamaan berlaku jika salah satu himpunan memuat himpunan
           lainnya, tetapi hal itu bukan syarat ``hanya jika'', sebagaimana
           ditunjukkan oleh contoh berikut di \( \Re^3 \).
           \begin{equation*}
             S=\set{\colvec[r]{1 \\ 0 \\ 0},\colvec[r]{0 \\ 1 \\ 0}},
             \quad
             T=\set{\colvec[r]{1 \\ 0 \\ 0},\colvec[r]{0 \\ 0 \\ 1}}
           \end{equation*}
         \partsitem Tidak pernah, karena rentang komplemen merupakan subruang,
           sedangkan komplemen rentang bukan merupakan subruang (komplemen itu
           tidak memuat vektor nol).
      \end{exparts}  
     \end{answer}
  % \item 
  %   Reprove \nearbylemma{le:SpanIsASubsp} without doing the
  %   empty set separately.
  %   \begin{answer}
  %     Call the subset \( S \).
  %     By \nearbylemma{th:SubspIffClosed},
  %     we need to check that 
  %     \( \spanof{S} \) is closed under linear combinations.
  %     If \( c_1\vec{s}_1+\dots+c_n\vec{s}_n,
  %       c_{n+1}\vec{s}_{n+1}+\dots+c_m\vec{s}_m\in\spanof{S} \) then
  %     for any \( p,r\in\Re \) we have
  %     \begin{multline*}
  %       p\cdot(c_1\vec{s}_1+\cdots+c_n\vec{s}_n)+
  %            r\cdot(c_{n+1}\vec{s}_{n+1}+\cdots+c_m\vec{s}_m)        \\
  %       =
  %       pc_1\vec{s}_1+\cdots+pc_n\vec{s}_n
  %         +rc_{n+1}\vec{s}_{n+1}+\cdots+rc_m\vec{s}_m
  %     \end{multline*}
  %     which is an element of \( \spanof{S} \).
  %     % (\textit{Remark.}
  %     % If the set $S$ is empty, then that
  %     % `if \ldots\ then \ldots' statement is vacuously true.)  
  %   \end{answer}
  \item Temukan suatu struktur yang tertutup terhadap kombinasi linear, tetapi
    bukan merupakan ruang vektor.
    % (\textit{Remark.} 
    % This is a bit of a trick question.)
    \begin{answer}
      Agar hal ini terjadi, salah satu syarat yang membuat operasi penjumlahan
      dan perkalian skalar berperilaku sebagaimana mestinya harus dilanggar.
      Pertimbangkan \( \Re^2 \) dengan operasi-operasi berikut.
      \begin{equation*}
        \colvec{x_1 \\ y_1}
        +\colvec{x_2 \\ y_2}
        =\colvec[r]{0 \\ 0}
        \qquad
        r\colvec{x \\ y}
        =\colvec[r]{0 \\ 0}
      \end{equation*}
      Himpunan $\Re^2$ tertutup terhadap operasi-operasi ini.
      Akan tetapi, himpunan tersebut bukan ruang vektor.
      \begin{equation*}
        1\cdot\colvec[r]{1 \\ 1}
        \neq\colvec[r]{1 \\ 1}
      \end{equation*}  
    \end{answer}
\index{subruang|)}
\index{ruang vektor|)}
\end{exercises}
